\documentclass[12pt,a4paper,openright,twoside]{article}
\usepackage{epsfig}
\usepackage[english]{babel}
\usepackage[utf8x]{inputenc}
\usepackage[T1]{fontenc}

\parskip=5pt
\parindent=15pt
\usepackage[margin=0.85in]{geometry}

%% Useful packages
\usepackage{amsmath}
\usepackage{amsthm}
\usepackage{amsfonts}
\usepackage{mathrsfs}
\usepackage{graphicx}
\usepackage{verbatim} 
\usepackage[algoruled,noline,longend]{algorithm2e}
\usepackage[colorinlistoftodos]{todonotes}
\DeclareMathOperator*{\E}{\mathbb{E}}
\DeclareMathOperator*{\R}{\mathbb{R}}
\DeclareMathOperator*{\N}{\mathbb{N}}
\DeclareMathOperator*{\Sspace}{\mathcal{S}}
\DeclareMathOperator*{\A}{\mathcal{A}}
\DeclareMathOperator*{\T}{\mathcal{T}}
\DeclareMathOperator*{\Z}{\mathcal{Z}}
\DeclareMathOperator{\Ppi}{\textbf{P}^\pi}
\DeclareMathOperator*{\Y}{\hat{\mathcal{Y}}_{\Hat{\gamma}}}

\numberwithin{equation}{section}
\newtheorem{theorem}{Theorem}[section]
\newtheorem{prop}[theorem]{Proposition}
\newtheorem{corol}[theorem]{Corollary}
\newtheorem{lemma}[theorem]{Lemma}
\newtheorem*{theorem*}{Theorem}
\newtheorem*{prop*}{Proposition}
\newtheorem*{corol*}{Corollary}

\theoremstyle{definition}
\newtheorem{definition}{Definition}
\newtheorem*{remark}{Remark}
\newtheorem*{notation}{Notation}

\theoremstyle{remark}
\newtheorem*{sol}{Resolution}

\theoremstyle{plain}
\newtheorem{mixtureProperty}{Property}

\def\qed{\hfill $\square$}

\graphicspath{ {./logos/} }

% --------------------------------------------------
\usepackage{fancyhdr}
\usepackage{emptypage}

\fancyhead[RO]{\small{\rightmark}}
\fancyhead[LE]{\small{\nouppercase{\leftmark}}}
\fancyhead[RE]{\small{Master Thesis}}
%\fancyhead[RE]{\small{Trabajo final de grado} \quad \quad \quad  \begin{picture}(0,0) \put(-32.84807,-3){\includegraphics[width=13mm]{./ub}} \end{picture} } 
\fancyhead[LO]{ \begin{picture}(0,0) \put(0,-3){\includegraphics[width=8mm]{./upc}} \end{picture} \quad \quad \begin{picture}(0,0) \put(0,-3){\includegraphics[width=13mm]{./ub}} \end{picture} \quad \quad \quad \begin{picture}(0,0) \put(0,-3){\includegraphics[width=9mm]{./urv}} \end{picture} \quad \quad  \quad\small{MAI} }
%\fancyhead[LO]{\begin{picture}(0,0) \put(0,-3){\includegraphics[width=13mm]{./ub}} \end{picture} \quad \quad \quad \small{Universitat de Barcelona} }
%\fancyhead[CE,CO]{\begin{picture}(0,0) \put(-17.01,-4.2){\includegraphics[width=12mm]{./ub}} \end{picture}} %foto de la ub en medio del encabezado
\fancyfoot[CO,CE]{}
\fancyfoot[RO,LE]{\thepage}

\fancypagestyle{plain}{%
\fancyhf{} % clear all header and footer fields
\fancyfoot[RO,LE]{\thepage} % except the center
\renewcommand{\headrulewidth}{0pt}
\renewcommand{\footrulewidth}{0pt}}

\setcounter{secnumdepth}{2}
\setcounter{tocdepth}{2}
\usepackage{hyperref}

%--------------------------------------------------

\begin{document}


\thispagestyle{empty}

Title


\newpage
\pagenumbering{roman} 

\newpage 


\section*{Aknowledgements}

\newpage


\section*{Abstract}

Temporal difference learning algorithms classically learn value functions (sum of discounted rewards) by solving for the fixed point that relates the value of one state with the expected value of the following states. Other TD learning algorithms have been used to learn \textit{multiplicative value functions}, which also have a fixed point solution. An example of such a function is the ratio between stationary distributions of two policies. It has been shown that learning of these multiplicative value functions suffers from higher variance and looser convergence guarantees. One potential approach to solving this issues would be to learn in log space, turning the multiplicative value function into an additive one. Unfortunately, it is easy to show with Jensen's inequality that the exponential of the log fixed point does not correspond to the multiplicative fixed point.

Distributional RL has been used to learn the return distribution (return is the sum of future discounted rewards). Although it has been shown that it mainly plays an auxiliary task role for representation learning, we propose that distributional fixed points could play a much more fundamental role to learning non additive value functions. In the case of multiplicative value functions, learning the additive fixed point in log space and exponentiating leads to the correct solution.

In particular, we will study this approach in the case of the learning process of the Covariate Shift ratio, which defines a multiplicative value function.





\newpage

%\pagenumbering{roman} 
\tableofcontents

\newpage


\pagestyle{fancy}
\renewcommand{\sectionmark}[1]{\markboth{\thesection.\quad #1}{}}
\renewcommand{\subsectionmark}[1]{\markright{\thesubsection.\quad #1}}

\pagenumbering{arabic} 
\setcounter{page}{1}

\thispagestyle{plain}
\section{Introduction}

- Reinforcement Learning

- Distributional Reinforcement Learning

- Off-Policy Learning

- Covariate Shift Ratio


%\subsection*{Project}


\subsection*{Motivation and Goals}

- Relevance of logarithmic Distributional RL

- Dessign of Distributional Covariate Shift Ratio

- Evaluation of the logarithmich approach within the Distributional CSR setting


\subsection*{Structure of the Report}




\newpage


\thispagestyle{plain}
\section{Distributional Reinforcement Learning}

TODO: INTRO

Initial paper\cite{DRL}:
In this paper we argue for the fundamental impor- tance of the value distribution: the distribution of the random return received by a reinforcement learning agent. This is in contrast to the com- mon approach to reinforcement learning which models the expectation of this return, or value. Although there is an established body of liter- ature studying the value distribution, thus far it has always been used for a specific purpose such as implementing risk-aware behaviour. We begin with theoretical results in both the policy eval- uation and control settings, exposing a signifi- cant distributional instability in the latter. We then use the distributional perspective to design a new algorithm which applies Bellman’s equa- tion to the learning of approximate value distri- butions. We evaluate our algorithm using the suite of games from the Arcade Learning En- vironment. We obtain both state-of-the-art re- sults and anecdotal evidence demonstrating the importance of the value distribution in approxi- mate reinforcement learning. Finally, we com- bine theoretical and empirical evidence to high- light the ways in which the value distribution im- pacts learning in the approximate setting.

Although the distributional perspective is almost as old as Bellman’s equation itself (Jaquette, 1973; Sobel, 1982; White, 1988), in reinforcement learning it has thus far been subordinated to specific purposes: to model parametric un- certainty (Dearden et al., 1998), to design risk-sensitive al- gorithms (Morimura et al., 2010b;a), or for theoretical anal- ysis (Azar et al., 2012; Lattimore  Hutter, 2012). By con- trast, we believe the value distribution has a central role to play in reinforcement learning.

Contraction of the policy evaluation Bellman operator. Basing ourselves on results by Ro sler (1992) we show that, for a fixed policy, the Bellman operator over value distribu- tions is a contraction in a maximal form of the Wasserstein (also called Kantorovich or Mallows) metric. Our partic- ular choice of metric matters: the same operator is not a contraction in total variation, Kullback-Leibler divergence, or Kolmogorov distance.

Instability in the control setting. We will demonstrate an instability in the distributional version of Bellman’s opti- mality equation, in contrast to the policy evaluation case. Specifically, although the optimality operator is a contrac- tion in expected value (matching the usual optimality re- sult), it is not a contraction in any metric over distributions. These results provide evidence in favour of learning algo- rithms that model the effects of nonstationary policies.

Better approximations. From an algorithmic standpoint, there are many benefits to learning an approximate distribu- tion rather than its approximate expectation. The distribu- tional Bellman operator preserves multimodality in value distributions, which we believe leads to more stable learn- ing. Approximating the full distribution also mitigates the effects of learning from a nonstationary policy. As a whole, we argue that this approach makes approximate reinforce- ment learning significantly better behaved.

From a supervised learning perspective, learning the full value distribution might seem obvious: why restrict our- selves to the mean? The main distinction, of course, is that in our setting there are no given targets. Instead, we use Bellman’s equation to make the learning process tractable; we must, as Sutton and Barto (1998) put it, “learn a guess from a guess”. It is our belief that this guesswork ultimately carries more benefits than costs.

WASSERSTEIN METRIC! THEORY DEVELOPED + The Cramer Distance as a Solution to Biased Wasserstein Gradients

INTERESTING APPENDIX ON RELATED WORK!!

It is surprising that, when we use a policy which aims to maximize expected return, we should see any difference in performance. The distinction we wish to make is that learning distributions matters in the presence of approxi- mation. We now outline some possible reasons.

Reduced chattering. Our results from Section 3.4 high- lighted a significant instability in the Bellman optimal- ity operator. When combined with function approxima- tion, this instability may prevent the policy from converg- ing, what Gordon (1995) called chattering. We believe the gradient-based categorical algorithm is able to mitigate these effects by effectively averaging the different distributions, similar to conservative policy iteration (Kakade AND Langford, 2002). While the chattering persists, it is inte- grated to the approximate solution.

State aliasing. Even in a deterministic environment, state aliasing may result in effective stochasticity. McCallum (1995), for example, showed the importance of coupling representation learning with policy learning in partially ob- servable domains. We saw an example of state aliasing in PONG, where the agent could not exactly predict the re- ward timing. Again, by explicitly modelling the resulting distribution we provide a more stable learning target.

A richer set of predictions. A recurring theme in artificial intelligence is the idea of an agent learning from a mul- titude of predictions (Caruana 1997; Utgoff and Stracuzzi 2002; Sutton et al. 2011; Jaderberg et al. 2017). The dis- tributional approach naturally provides us with a rich set of auxiliary predictions, namely: the probability that the return will take on a particular value. Unlike previously proposed approaches, however, the accuracy of these pre- dictions is tightly coupled with the agent’s performance.

Framework for inductive bias. The distributional perspective on reinforcement learning allows a more natural framework within which we can impose assumptions about the domain or the learning problem itself. In this work we used distributions with support bounded in [VMIN,VMAX]. Treating this support as a hyperparameter allows us to change the optimization problem by treating all extremal returns (e.g. greater than VMAX ) as equivalent. Surprisingly, a similar value clipping in DQN significantly degrades per- formance in most games. To take another example: in- terpreting the discount factor as a proper probability, as some authors have argued, leads to a different algorithm.


Well-behaved optimization. It is well-accepted that the KL divergence between categorical distributions is a rea- sonably easy loss to minimize. This may explain some of our empirical performance. Yet early experiments with al- ternative losses, such as KL divergence between continu- ous densities, were not fruitful, in part because the KL di- vergence is insensitive to the values of its outcomes. A closer minimization of the Wasserstein metric should yield even better results than what we presented here.

An Analysis of Categorical Distributional Reinforcement Learning

Distributional Reinforcement Learning with Quantile Regression

Implicit Quantile Networks for Distributional Reinforcement Learning

DISTRIBUTED DISTRIBUTIONAL DETERMINISTIC POLICY GRADIENTS

A Comparative Analysis of Expected and Distributional Reinforcement Learning

Statistics and Samples in Distributional Reinforcement Learning:




\subsection{Reinforcement Learning Setting}

We consider an agent interacting with an environment in the standard setting\cite{rlformulation}: at each step $t$, the agent selects an action $a_t$ based on its current state $s_t$, to which the environment responds with a reward $r_t$ and then moves to the next state $s_{t+1}$. We model this interaction as a time-homogeneous Markov Decision Process $(\Sspace, \A, r, P, \gamma)$, where
\begin{itemize}
    \item $\Sspace$ and $\A$ are the state and action spaces, respectively, we assume that both are finite, with $n:=|\Sspace|$;
    \item $P$ is the transition kernel, $s_{t+1} \sim P(\cdot | s_t,a_t)$; the Markov assumption states that \\ $P(s_{t+1} | s_t, a_t, s_{t-1}, a_{t-1}, ...) = P(s_{t+1} | s_t, a_t)$;
    \item $r(s,a)$ represents the immediate reward given by the environment after taking action $a$ being in state $s$. These rewards are considered to be sampled from the reward function $R(s,a)$, i.e. $r_t \sim R(s_t,a_t)$;
    \item $\gamma$ is the discount factor
\end{itemize}

A policy $\pi$ maps each state to a probability distribution over the action space, $a_t \sim \pi(\cdot | s_t)$. In addition, we combine the policy $\pi$ and transition function $P$ into a state-to-state transition function $P_\pi \in \R^{n \times n}$ , whose entries are
\begin{equation} \label{transitionFunction}
    P_{\pi} (s'|s) := Prob_\pi (s_{t+1}=s' | s_t = s) = \sum_{a\in\A} \pi(a|s)P(s'|s,a)
\end{equation}
In particular, powers of $P_\pi$ represent the transition function across different time-steps.

Given $\pi$, the action value function is defined as the expected sum of discounted rewards from a state-action pair by following the policy:
\begin{equation} \label{AV}
Q^\pi (s,a) = \mathbb{E}_\pi \left[ \sum_{t=0}^\infty \gamma^t r(s_t,a_t) \bigg| s_0=s, a_0=a \right] 
\end{equation}
The Bellman's equation can be obtained from this expression:
\begin{equation} \label{BE}
\begin{split}
    Q^\pi (s,a) &= \mathbb{E} \left[ r(s,a) \right] + \gamma \E_\pi \left[ \sum_{t=0}^\infty \gamma^t r(s_{t+1},a_{t+1}) \bigg| s_0=s, a_0=a \right] \\
    &= \mathbb{E} \left[ r(s,a) \right] + \gamma \sum_{s'} P(s'|s,a) \left( \E_\pi \left[ \sum_{t=0}^\infty \gamma^t r(s_{t+1},a_{t+1}) \bigg| s_1=s' \right] \right) \\
    &= \mathbb{E} \left[ r(s,a) \right] + \\
    & \quad \gamma \sum_{s'} P(s'|s,a) \sum_{a'} \pi(a'|s')\left( \E_\pi \left[ \sum_{t=0}^\infty \gamma^t r(s_{t+1},a_{t+1}) \bigg| s_1=s', a_1=a' \right] \right)\\
    %&= \mathbb{E} \left[ r(s,a) \right] + \gamma \sum_{s'} P(s'|s,a) \sum_{a'} \pi(a'|s') Q(s',a') \\
    &= \mathbb{E} \left[ r(s,a) \right] + \gamma \E_{s'\sim P(\cdot | s,a), a' \sim \pi(\cdot | s')} \bigg[Q(s',a') \bigg] \\
\end{split}
\end{equation}


Analogously for the state value function (considering $r_\pi(s) := \mathbb{E}_{a\sim\pi(\cdot | s)} \left[ r(s,a) \right]$):
\begin{equation} \label{SV}
\begin{split}
    V^\pi (s) &= \mathbb{E}_\pi \left[ \sum_{t=0}^\infty \gamma^t r(s_t,a_t) \bigg| s_0=s\right] \\
    &= r_\pi(s) + \gamma \sum_a \pi(a | s) \sum_{s'} P(s'|s,a) \left( \E_\pi \left[ \sum_{t=0}^\infty \gamma^t r(s_{t+1},a_{t+1}) \bigg| s_1=s' \right] \right) \\
    &= r_\pi(s) + \gamma \E_{s'\sim P_{\pi}(\cdot | s)} \bigg[V^\pi(s') \bigg]
\end{split}
\end{equation}

In the context of \textit{policy evaluation}, and considering state-action value functions $Q^\pi$ as vectors in $\R^{\Sspace \times \A}$, the \textit{Bellman operator} $\T^{\pi} : \R^{\Sspace \times \A} \rightarrow \R^{\Sspace \times \A}$ is defined as
\begin{equation} \label{BellmanOp}
	{\T}^\pi Q(s,a) :=  \mathbb{E} \left[ r(s,a) \right] + \gamma \E_{s'\sim P(\cdot | s,a), a' \sim \pi(\cdot | s')} \bigg[Q(s',a') \bigg]
\end{equation}
Such operator is useful to describe the expected behaviour of popular learning algorithms (e.g. Q-learning) and satisfies some interesting and desirable properties\cite{rlformulation}:
\begin{itemize}
	\item It is a contraction mapping (DEFINE IT, OR APPENDIX)
	\item The process $Q_t = \T^\pi Q_{t-1}$, for some initial value $Q_0$, converges exponentially to $Q^\pi$ as $t \rightarrow \infty$.
\end{itemize}


Regarding the \textit{control setting}, where the goal is to improve the current policy $\pi$, a \textit{Bellman optimality operator} $\T^* : \R^{\Sspace \times \A} \rightarrow \R^{\Sspace \times \A}$ is defined as follows:
\begin{equation} \label{BellmanOptimalityOp}
	{\T}^* Q(s,a) :=  \mathbb{E} \left[ r(s,a) \right] + \gamma \E_{s'\sim P(\cdot | s,a)} \bigg[ \max_{a' \in \A}  Q(s',a') \bigg]
\end{equation}
As it was shown in \cite{rlformulation}, $\T^*$ is also a contraction mapping, and the expectation of $\{ Q_t \}$, $Q_t = \T^* Q_{t-1}$, converges exponentially to a fixed point (in this case, the optimal value function $Q^*$).

\subsection{Towards a Distributional Reinforcement Learning}

TODO: INTRO

In \cite{DRL} a distributional Bellman equation is defined:
\begin{equation} \label{distBE}
    Z^{\pi}(s,a) \stackrel{D}{=} R(s,a) + \gamma Z^{\pi}(S',A')
\end{equation}
where
\begin{itemize}
    \item $Z^\pi: \Sspace \times \A \rightarrow \Z$ is the \textit{value distribution}, a mapping from state-action pairs to distributions over returns by following policy $\pi$. Its expectation is the value $Q^\pi$
    	\begin{equation} \label{QexpZ}
        		Q^\pi(s,a) := \mathbb{E} [Z^\pi(s,a)]
    	\end{equation}
    	We will also called it the \textit{random return}.
    \item $(S',A') \in \Sspace \times \A$ is the next state-action random variable: $S'\sim P(\cdot | s,a)$, $A' \sim \pi(\cdot | S')$
    \item $R(s,a) \in \Z$ is the random reward, or equivalently the reward function. Note that now we are dealing with it as an explicit random variable.
    \item $Z^\pi(S',A') \in \Z$ is the random return over the random next state-action following $\pi$. This notation implies that all possible next state-action pairs need to be considered as to generate this return distribution. Thus, $Z^\pi(S',A')$ may be seen as a mixture distribution of the distributions $Z^\pi(s',a')$ where $s'$ and $a'$ are sampled from $(S',A')$:
    \begin{equation}
        f_{Z^\pi(S',A')} (z) = \sum_{s'} P(s' | s,a) \sum_{a'} \pi(a'|s') f_{Z^\pi(s',a')}(z)
    \end{equation}
    The expected value, using \ref{QexpZ}, can be then expressed as:
    \begin{equation} \label{expNextRandomReturn}
    \begin{split}
        \E[Z^\pi(S',A')] &=  \int_{-\infty}^{\infty} z \sum_{s'} P(s' | s,a) \sum_{a'} \pi(a'|s') f_{Z^\pi(s',a')}(z) dz \\
        &=  \sum_{s'} P(s' | s,a) \sum_{a'} \pi(a'|s') \int_{-\infty}^{\infty} z f_{Z^\pi(s',a')}(z) dz \\
        &=  \sum_{s'} P(s' | s,a) \sum_{a'} \pi(a'|s') \E[Z^\pi(s',a')] \\
        &= \E_{s'\sim P(\cdot | s,a), a' \sim \pi(\cdot | s')} \bigg[Q^\pi(s',a') \bigg] 
    \end{split}
    \end{equation}
\end{itemize}

\begin{remark}
A distributional equation $U \stackrel{D}{:=} V$ indicates that the random variable $U$ is distributed according to the same law as $V$.	
\end{remark}

Note that we can easily recover the classical Bellman's equation \ref{BE} for the action value Q by using \ref{QexpZ} and \ref{expNextRandomReturn} when taking the expected value over its distributional version \ref{distBE}:
\begin{equation}
\begin{split}
    Q^\pi(s,a) &= \mathbb{E} [Z^\pi(s,a)] \\
    &= \E[R(s,a)] + \gamma \E[Z^\pi(S',A')] \\
    &= \mathbb{E} [r(s,a)] + \gamma \E_{s'\sim P(\cdot | s,a), a' \sim \pi(\cdot | s')} \bigg[Q^\pi(s',a') \bigg] 
\end{split}
\end{equation}

Finally, let's try to find out what actually the random return $Z$ represents by expanding its density function: 
\begin{equation} \label{meaningZ}
\begin{split}
    f_{Z^\pi(s,a)}(z) &= f_{R(s,a) + \gamma Z^\pi(S',A')}(z)  \\
    &=  \sum_{s'} P(s' | s,a) \sum_{a'} \pi(a'|s') f_{R(s,a) + \gamma Z^\pi(s',a')}(z) \\
    &= \sum_{s'} P(s' | s,a) \sum_{a'} \pi(a'|s') f_{R(s,a) + \gamma (R(s',a') + \gamma Z^\pi(S'',A''))}(z) \\
    &= \sum_{s'} P(s' | s,a) \sum_{a'} \pi(a'|s') \sum_{s''} P(s'' | s',a') \sum_{a''} \pi(a''|s'') \\
    & \quad \quad f_{R(s,a) + \gamma R(s',a') + \gamma^2 Z^\pi(s'',a'')} (z) \\
    &= \sum_{s_{1}} P(s_{1} | s_{0},a_{0}) \sum_{a_{1}} \pi(a_{1}|s_{1})  \cdots \sum_{s_{t}} P(s_{t} | s_{t-1},a_{t-1}) \sum_{a_{t}} \pi(a_{t}|s_{t}) \\
    & \quad \quad f_{\sum_{i=0}^{t} \gamma^t R(s_{t},a_{t})} (z)
\end{split}
\end{equation}
where we have repeatedly used property \ref{mixtureProperty2} of mixture distributions. In addition, note that assuming independence between the random reward $R$ of a certain state-action pair and the return distribution $Z^\pi$ of the possible next state-action (which is NOT TRUE in general), we can rewrite it in terms of convolutions as
\begin{equation}
\begin{split}
    f_{Z^\pi(s,a)}(z) &= \sum_{s_{1}} P(s_{1} | s_{0},a_{0}) \sum_{a_{1}} \pi(a_{1}|s_{1})  \cdots \sum_{s_{t}} P(s_{t} | s_{t-1},a_{t-1}) \sum_{a_{t}} \pi(a_{t}|s_{t}) \\
    & \quad \quad  \left( f_{R(s_0,a_0)} * f_{\gamma R(s_1,a_1)} *  \cdots * f_{\gamma^t R(s_{t},a_{t})} \right) (z)
\end{split}
\end{equation}

According to \ref{meaningZ}, $Z^{\pi}(s,a)$ can be interpreted as a convex combination of the sum of discounted reward distributions of all possible agent trajectories starting at the state-action $(a,s)$ and following policy $\pi$ from then on, each weight corresponding to the probability of that precise trajectory. It is important to note that $Z^\pi$ encodes the intrinsic randomness of the agent's interactions with its environment; we should avoid considering it as a measure of uncertainty about the environment itself.

Moving to the \textit{policy evakluation} setting, now we are interested in studying the behaviour of a distributional version of the policy evaluation operator $\T^\pi$. First we define the transition operator $P^\pi : \Z \rightarrow \Z$ 
\begin{equation} \label{distTransitionOp}
	P^\pi Z(s,a) \stackrel{D}{:=} Z(S',A')
\end{equation}
Given this, the \textit{distributional Bellman operator} $\T : \Z \rightarrow \Z$ is defined as
\begin{equation} \label{distBellmanOp}
	{\T}^\pi Z(s,a) \stackrel{D}{:=} R(s,a) + \gamma P^{\pi} Z(s,a)
\end{equation}
We emphasize that three sources of randomness are involved in the compound distribution ${\T}^\pi Z$, i.e.
\begin{enumerate}
	\item Randomness in the reward $R$,
	\item Randomness in the transition $P^\pi$, and
	\item Randomness in the next state-value distribution $Z(S', A')$,
\end{enumerate} 
which together make this distributional Bellman operator fundamentally different to the expected value-cased one (Equation \ref{BellmanOp}). Authors in \cite{DRL} demonstrate, under the assumption that these three random quantities are independent, that ${\T}^\pi$ is a contraction mapping whose unique fixed point is the value distribution $\Z^\pi$.


However, more difficulties arise in the \textit{control setting} when dealing with this distributional perspective; as stated in \cite{DRL}, while every optimal policy attain the same value $Q^*$ in the expected-valued case, there might be many optimal value distribiutions. Considering the set $\Pi^*$ of optimal policies, an optimal value distribution is precisely characterized as follows
\begin{definition}
	An optimal value distribution is the value distribution of an optimal policy. The set of optimal value distributions is $\Z^* := \{ Z^{\pi^*} : \pi^* \in \Pi^* \}$.
\end{definition} 
Note that an optimal value distribution must match the full distribution of returns under some optimal policy, so that not all value distributions with expectation $Q^*$ are optimal. A \textit{distributional Bellman optimality operator} $\T$ is defined as well, but it doesn't behave as well as the policy evaluation operators: it is not a contraction in any usual metric between distributions, and its convergence to the set of optimal value distributions is weak.



\subsection{Approximation Framework in the Distributional Setting}

%Let $\Ppi (s,a)$ be the distribution of the random variable $Z^\pi (s,a)$. The process of learning $\Ppi$ 
We should we aware that the full computation of the distributional Bellman operators on return distribution functions is generally either impossible (as we typically do not have access to the MDP dynamics, but to merely sample transitions) or infeasible (since the value distribution cannot be stored exactly in the general case).

This lead us to the design of several key approximations which are required to implement practical and scalable distributional RL algorithms\cite{analysisC51}:

\subsubsection{Distribution Parametrisation}
Due to the fact that the full space of probability distributions, $\mathscr{P}(\R)$, cannot be algorithmically encoded with a finite number of parameters, we need to approximate the distribution throughout a parametric family $\mathcal{P} \in \mathscr{P}(\R)$.

\subsubsection{Stochastic Bellman Operators}
In order to evaluate the distributional Bellman operator $\T^\pi$, all possible next state-action-reward combinations should be taken into account. As in the expected valued case, the usual way of overcoming this practical limitation is by learning through transition samples $(s,a,r,s',a')$ of the MDP. Hence, we can define a \textit{stochastic distributional Bellman operator} $\widehat{\T}^\pi$ adapted to the randomness of these transitions, which defines a random measure whose behaviour is equal in expectation to the true Bellman operator $\T^\pi$. %In our theoretical development, for the sake of simplicity, we will consider the simpler setting where stochastic operators are not needed.

\subsubsection{Projection of Bellman Target Distribution}
Another problem usually arises after computing the stochastic operator $\widehat{\T}^\pi$ over a value distribution $Z(s,a)$: the new distribution may no longer lie in the selected parametric family $\mathcal{P}$. If this is the case, we further need to apply a \textit{projection operator} $\Pi : \mathscr{P}(\R) \rightarrow \mathcal{P}$ so as to map $\widehat{\T}^\pi Z(s,a)$ into the proper parametric family.

\subsubsection{Gradient Updates}
 Finally, having computed a stochastic approximation $\widehat{Z}_k (s_k,a_k) = \Pi \widehat{\T}^\pi Z_k(s_k,a_k)$ to the full target distribution, we still have to define how to compute the next iterate $Z_{k+1}$. For that, the use of gradient updates seems to be appropiate\cite{rlformulation}, as it helps dissipate some noise introduced in the target by the stochastic approximation. A key aspect here, however, is the not-straightforward correspondence between the considered loss and the previously selected projection; convergence and good behaviour of the resulting algorithm have only be proven when both of them rely on the same norm-induced geometry\cite{analysisC51,rlformulation}

\subsection{Categorical Distributional Reinforcement Learning}

Although the distributional perspective is almost as old as Bellman's equations\cite{firstDistributional}, it was not until the recent introduction of the Categorical Distributional Reinforcement\cite{DRL} (CDRL) that it has become a central role within reinforcement learning. Their algorithm, called C51, was able to obtain state-of-the-art results in the Arcade Learning Environment\cite{ALE} (ALE), outperforming the top expected-valued solutions by then.

The name 'categorical' comes from the distribution parametrisation that is used, which consists in the parametric family of categorical distributions over some fixed set of equally-spaced supports $z_1 < \cdots < z_K$:
\begin{equation}
	\mathcal{P} = \left\{ \sum_{i=1}^{K} p_i \delta_{z_i} \Big| p_1, \dots , p_K \geq 0, \sum_{k=1}^{K} p_k = 1 \right\}
\end{equation}
Apart from that, the algorithm can be fully characterized by the rest of approximations implemented within the distributional RL framework described in the previous subsection:
\begin{itemize}
	\item So as to learn from sampling transitions, we move to the \textit{stochastic distributional Bellman operator} $\widehat{\T}^\pi$ and the \textit{stochastic distributional Bellman optimality operator} $\widehat{\T}^*$. Given a sampled transition $(s_t,a_t,r,s_{t+1},a_{t+1})$, these stochastic operators basically transform the supports of the distributions by an affine shift map $f_{r,\gamma} : \R \rightarrow \R$, defined by $f_{r,\gamma}(z) = r + \gamma z$; in our notation, 
		$$\widehat{\T} Z_t(s_t,a_t) = (f_{r,\gamma})_\# Z_t(s_{t+1},a_{t+1})$$
		with $a_{t+1}$ either being selected by sampling the policy $\pi(\cdot | s_{t+1})$  (categorical policy evaluation, i.e. $\widehat{\T}=\widehat{\T}^\pi$), or being the action with the highest estimated expected returns (categorical Q-learning, $\widehat{\T}=\widehat{\T}^*$).
		
	\item CDRL applies the heuristic projection operator $\Pi_C$ after computing the stochastic Bellman operators in order to recover a distribution within the selected parametric family $P$. This projection is defined for single Dirac measures as
	\begin{equation} \label{heuristicProjectionOp}
		\Pi_C (\delta_y) = \left\{
		\begin{array}{ll}
			\delta_{z_1} & y\leq z_1 \\
			\\ \frac{z_{i+1} -y}{z_{i+1}-z_i} \delta_{z_i} + \frac{y-z_i}{z_{i+1}-z_i} \delta_{z_{i+1}} & z_i < y < z_{i+1} \\
			\\ \delta_{z_K} & y>z_K
		\end{array}
		\right.
	\end{equation}
	and can be easily extended to finite mixtures of Dirac measures: 
	$$ \Pi_C \left( \sum_{i=1}^{N} p_i \delta_{y_i} \right) = \sum_{i=1}^{N} p_i \Pi_C ( \delta_{y_i} )	$$
	This projection step provides us with the target $\widehat{Z}_t (s_t,a_t) = \Pi_{C} \widehat{\T} Z_t(s_t,a_t)$. 
	
	In addition, we will also denote by $\Pi_C$ the operator $\mathscr{P}(\R)^{\Sspace \times \A} \rightarrow \mathcal{P}^{\Sspace \times \A}$ that applies this projection to each distribution in the return distribution function.
	
	\item Finally, the original C51 performs a single step of gradient descent on the Kullback-Leibler divergence (see Appendix \ref{metrics} for the definition of this metric) of the prediction $Z_t(s_t,a_t)$ from the target $\widehat{Z}_t (s_t,a_t)$:
	$$ \textnormal{KL} \left( \widehat{Z}_t (s_t,a_t) || Z_t (s_t,a_t) \right) $$
	with respect to the parameters of $Z_t (s_t,a_t)$; this gradient is used to generate the new estimate $Z_{t+1}(s_t,a_t) = \sum_{k=1}^{K} p_{t+1,k}(s_t,a_t) \delta_{z_k}$.
	
\end{itemize}

Pseudo-Algorithm \ref{CDRL} (provided by \cite{analysisC51}) synthesises the steps performed by CDRL.

\setlength{\algomargin}{1.5em}
\SetAlCapHSkip{0em}
\begin{algorithm}[H]
\caption{Categorical Distributional Reinforcement Learning} \label{CDRL}
\DontPrintSemicolon
\SetKwComment{CustomComment}{\#}{}
\SetKwInput{KwIn}{\hspace{-1.5em} Input}
\SetKwInput{KwOut}{\hspace{-1.5em} Output}
\SetKwInOut{Require}{\hspace{-1.5em} Require}
\Require{$Z_t(s,a) = \sum_{k=1}^K p_{t,k}(s,a) \delta_{z_k}$ for each $(s,a)$}
\KwIn{A sample transition $(s_{t},a_{t}, r_t, s_{t+1})$}
	\CustomComment{Compute distributional Bellman target}
    \uIf{\textnormal{Categorical Policy Evaluation}}{
    	$a^* \sim \pi(\cdot | s_{t+1})$ \;}
    \ElseIf{\textnormal{Categorical Q-learning}}{
    	$a^* \leftarrow \textnormal{arg max}_a Q(s_{t+1},a) $\;}
    $\widehat{Z}_{*}(s_t,a_t) \leftarrow (f_{r_t,\gamma})_{\#} Z_t(s_{t+1},a^*)$\;
    \CustomComment{Project target onto support}
    $\widehat{Z}_{t}(s_t,a_t) \leftarrow \Pi_C \widehat{Z}_{*}(s_t,a_t)$\;
    \CustomComment{Compute KL loss}
    \textnormal{Find gradient }$ \textnormal{KL} \left( \widehat{Z}_t (s_t,a_t) || Z_t (s_t,a_t) \right) $\; 
    \textnormal{Generate new estimate }$Z_{t+1}(s_t,a_t) = \sum_{k=1}^{K} p_{t+1,k}(s_t,a_t) \delta_{z_k}$ \;
\KwOut{Estimate $Z_{t+1}(s,a)$ for each $(s,a)$}

\end{algorithm}


Despite the good results obtained by C51, all attempts to prove its convergence have failed so far. One of its most controversial steps, at least from a theoretical point of view, is the use of the KL divergence while applying the heuristic projection $\Pi_{C}$; as shown in \cite{analysisC51}, using the Cramér distance instead allows a useful geometrical reinterpretation of $\Pi_{C}$ that lead us to the desired convergence results.

In fact, we recall that authors in \cite{DRL} demonstrates, in the context of policy evaluation setting, that applying the Bellman operator $T^{\pi}$ repeatedly to an initial return distribution function $Z_0$ guarantees convergence to the true set of return distributions $Z^\pi$ in the supremum-Wasserstein metric. They did so by showing that
\begin{lemma} \label{contractiondp}
	$\T^\pi : \Z \rightarrow \Z$ is a $\gamma$-contraction in $\bar{d}_p$
\end{lemma}

However, after the introduction of the parametrization $\mathcal{P}$ and the projection operator $\Pi_C$, it is important to note that the operator to be analyzed is not the Bellman operator $\T^\pi$ itself anymore, but its composition with the projection operator, i.e. $\Pi_C \T^\pi$. As stated in \cite{analysisC51}, this can make contractivity break under all Wasserstein distances but $\bar{d}_1$ :
\begin{lemma} \label{nocontractiondp}
	$\Pi_C \T^\pi : \Z \rightarrow \Z$ is in general not a contraction in $\bar{d}_p$, for $p>1$
\end{lemma}

In contrast to that, Cramér distance seems to induce a useful geometric structure on the space of probability measures, allowing us to establish contractivity of the combined operator $\Pi_C \T^\pi$ in a much more natural way. We dedicate the following two sections to analyze CDRL algorithms together with Cramér-based metrics in the case of both Tabular Representations and Linear Function Approximations, respectively, which are the theoretical frameworks where stable behaviour and convergence properties are proven so far.


\subsection{CDRL with Tabular Representation} \label{CDRLtabular}

Our first approach to study Cramér-based CDRL algorithms will be to consider the simplest reinforcement learning framework, where we are able to store an approximate parametrization distribution for each state-action pair: the so-called tabular case.

Again, we will consider the parametric family $\mathcal{P}$ of categorical distributions over some fixed support $\{z_1, \dots , z_K\}$, as well as the heuristic projection operator $\Pi_C$ defined in Equation \ref{heuristicProjectionOp} for mapping the backup distribution function $\T^\pi Z$ into $\mathcal{P}$. However, for the sake of simplicity and understandability, we will assume that no stochastic approximation is required. 

First of all, we present the result of \cite{analysisC51} that motivates the use of Cramér distance, which shows an interesting connection between this distance and the projection operator $\Pi_C$:
\begin{prop} \label{CramerProj}
	The Cramér metric $\ell_2$ endows a particular subset of $\mathscr{P}(\R)$ with a notion of orthogonal projection, and the orthogonal projection onto the subset $\mathcal{P}$ is exactly the heuristic projection $\Pi_C$. Consequently, $\Pi_C$ is a non-expansion with respect to $\ell_2$.
\end{prop}

The previous proposition directly lead us to the contractivity of the operator $\Pi_C \T^\pi$, which in turn ensures its convergence in the absence of stochastic approximation:
\begin{prop} \label{CramerBellmanContraction}
	The operator $\Pi_C \T^\pi$ is a $\sqrt{\gamma}$-contraction in $\bar{\ell}_2$. Further, there is a unique distribution function $Z_C \in \mathcal{P}^{\Sspace \times \A}$ to which the process $Z_{k+1} := \Pi_c \T^\pi Z_k$, given any initial distribution function $Z_0 \in \mathscr{P}(\R)^{\Sspace \times \A}$, converges in $\bar{\ell}_2$ as $k \rightarrow \infty$.
\end{prop}

Authors in \cite{analysisC51} also address the natural question of how the limiting distribution function $Z_C$ can differ from the true return distribution $Z^\pi$:
\begin{prop}
	Let $Z_C$ be the limiting return distribution of Proposition \ref{CramerBellmanContraction}. If for all state-action pairs their corresponding true return distribution $Z^\pi$ is supported on $[z_1,z_K]$, then 
	$$\bar{\ell}_2(Z_C, Z^\pi) \leq \frac{1}{1-\gamma} \max_{1\leq i \leq K} (z_{i+1}-z_i). $$
\end{prop}
Note that the true return distribution is gradually recovered as the fineness of the support increases. Hence, this result can be viewed as a way of quantifying the cost of using the parametrization $\mathcal{P}$ instead of fully non-parametric probability distributions. Finally, we also highlight the assumption that the support of the true return distributions lie on the selected support, since we don't always have access to the scale of rewards in every RL problem; for these circumstances, an analogous result can be found in Proposition 4 of \cite{analysisC51}.

\setlength{\algomargin}{1.5em}
\SetAlCapHSkip{0em}
\begin{algorithm}[H]
\caption{CDRL with Tabular Representation} \label{tabularCDRL}
\DontPrintSemicolon
\SetKwComment{CustomComment}{\#}{}
\SetKwInput{KwIn}{\hspace{-1.5em} Input}
\SetKwInput{KwOut}{\hspace{-1.5em} Output}
\SetKwInOut{Require}{\hspace{-1.5em} Require}
\Require{$Z_t(s,a)$ for each $(s,a)$}
\KwIn{A state-action pair $(s_{t},a_{t})$}
	\CustomComment{Compute distributional Bellman target}
    \uIf{\textnormal{Categorical Policy Evaluation}}{
    	$\widehat{Z}(s_t,a_t) \leftarrow \T^\pi Z_t(s_{t},a_{t})$\;}
    \ElseIf{\textnormal{Categorical Q-learning}}{
    	$\widehat{Z}(s_t,a_t) \leftarrow \T^* Z_t(s_{t},a_{t})$\;}
    \CustomComment{Project target onto support}
    $\widehat{Z}_{t}(s_t,a_t) \leftarrow \Pi_C \widehat{Z}(s_t,a_t)$\;
    \CustomComment{Update estimate}
    $Z_{t+1}(s_t,a_t) \leftarrow \widehat{Z}_{t}(s_t,a_t)$ \;
\KwOut{Estimate $Z_{t+1}(s,a)$ for each $(s,a)$}

\end{algorithm}

\subsection{CDRL with Linear Function Approximation}

In this second and last theoretical approach, we go a step beyond the tabular case by considering a framework where we make use of function approximators to compute the parametrized return distributions of each state-action pair, assuming a more realistic situation in which the state-action space is too large to store individual distributions.

We will restrict our analysis to the context of linear function approximation, whose convergence is already proven in \cite{DRLlinear}. To the best of our knowledge, no Distributional RL algorithm with a non-linear function approximator has been proven to converge so far, despite the many efforts in studying C51.

Regarding the rest of approximations and assumptions of the CDRL algorithms considered in this study, they are exactly the same as in the previous Section \ref{CDRLtabular}: we contemplate the parametric family $\mathcal{P}$ with fixed support $\{z_1, \dots, z_K\}$, the Cramér-based projection $\Pi_C$, and reject any stochastic approximation.

In our linear model, we now redefine the approximated return distributions $Z$ as mappings from states $s\in \Sspace$ to vectors defined by a linear combination of features:
\begin{equation}
	Z_{\Theta}(s) := \Theta^\top \phi(s)
\end{equation}
where $\phi(s)\in \R^m$ is the feature vector at state $s$ and $\Theta \in \R^{n\times K}$ represents the weight matrix. Therefore, we are assuming that the vector $Z_{\Theta}(s) \in \R^{K}$ is an estimation of a distribution over the support $\{z_1,\dots,z_K\}$, although it might have negative components and it is not necessarily normalized. 

\begin{notation}
	In vector notation, we will write $Z_{\Theta} = \Phi \Theta \in \R^{n\times K}$, where $\Phi \in \R^{n \times m}$ is the feature matrix, simply formed by all feature vectors.
\end{notation}


Before continuing, let's carefully review the three different spaces of distributions-like objects that we face in CDRL with linear function approximation:
\begin{itemize}
	\item On the one hand, we have the probability space where true return distributions lie; without any prior knowledge, we should simply consider the full space of distributions $Z^\pi \in \mathscr{P}(\R)$. Assuming that we know \textit{a priori} the scale of the rewards, we can reduce it to the space of distributions with support in a certain interval; in our case, it could be $Z^\pi \in \mathscr{P}([z_1,z_K])$.
	\item On the other hand, the selected parametric family $\mathcal{P}$ of categorical distributions over the support $\{z_1,\dots,z_K\}$; so far, we have expected our estimated return distributions $Z$ to belong to this space, $Z \in \mathcal{P}$.
	\item Finally, and in addition to the previous two, we now have a new vector space spanned by the features $\Phi \in \R^{n\times m}$, i.e. $\mathcal{L} := \{ \Phi\Theta : \Theta \in \R^{m\times K}\}$. This is, in fact, the crucial space in our development, as every approximated return distribution $\Z_{\Theta}$ lies in it. Ideally, if all $Z_\Theta$ were proper probability distributions over $\{z_1,\dots,z_K\}$, we would have $\mathcal{L} \subset \mathcal{P}$ (assuming some loss of expressiveness by our linear model).
\end{itemize}

%This can be seen in Pseudo-algorithm \ref{linearCDRL}, which defines the process of linear-based CDRL step by step. 
Let's review the process. First, our linear function model provides us with return distribution estimates 
$$Z_\Theta \in \mathcal{L}$$
 for each state $s \in \Sspace$. In order to get better approximations, CDRL algorithm applies the distributional Bellman operator, so we get
 $${\T}^\pi Z_{\Theta} \in \mathscr{P}([z_1,z_K]).$$
 Then we would need to apply a first projection operator $\Pi_1 : \mathscr{P}([z_1,z_K]) \rightarrow \mathcal{P}$ (like $\Pi_C$) so as to recover a distribution within our selected parametric family, so
 $$\Pi_1 {\T}^\pi Z_{\Theta} \in \mathcal{P}.$$
 However, for the algorithm to compare distributions and improve its linear representations, now we need a further projection onto the space $\mathcal{L}$, which defines all the probability distributions that our linear model can actually express. Thus, the idea is that a second projection operator $\Pi_{2} : \mathcal{P} \rightarrow \mathcal{L}$ is required so that
 $$\Pi_2 \Pi_1 {\T}^\pi Z_{\Theta} \in \mathcal{L}.$$
As shown in \cite{DRLlinear}, the (re)definition of both projections plays again a very important role for proving the convergence of the process $Z_{\Theta_{k+1}} = \Pi_2 \Pi_1 {\T}^\pi Z_{\Theta_{k}}$. 

So far, we demonstrated in the tabular case that the combined operator $\Pi_C \T^\pi$ is a contraction mapping in the supremum-Cramér metric $\bar{\ell}_2$; our task now is to achieve a similar result for the combination $ \Pi_2 \Pi_1 {\T}^\pi$. For doing so, authors in \cite{DRLlinear} introduce a \textit{generalized Cramér distance} which is able to deal with our return distribution estimates (we recall they are not necessarily probability distributions).

We denote denote by $C\in \R^{K\times K}$ the lower-triangular matrix of ones, and by $e\in\R^K$ the vector of ones. Let $e_K= (1/\sqrt{K}) e$ and $\Pi_{e_K^\top} = I_K - e_K e_K^\top$. For any two discrete value distributions $Z_1, Z_2$ over the fixed support $\{z_1,\dots,z_K\}$, the generalized Cramér distance is defined as
\begin{equation} \label{generalizedCramer}
	\ell_{\lambda}^2(Z_1,Z_2) = (Z_1 - Z_2)^\top \Pi_{e_K^\top} C C^\top \Pi_{e_K^\top} (Z_1 - Z_2) + \lambda ((Z_1 - Z_2)^\top e_K)^2 
\end{equation}
We will simplify the notation by denoting $C_\lambda = \Pi_{e_K^\top} C C^\top \Pi_{e_K^\top} + \lambda e_K e_K^\top $, so we can rewrite the distance definition as 
\begin{equation*}
	\ell_{\lambda}^2(Z_1,Z_2) = (Z_1 - Z_2)^\top C_\lambda (Z_1 - Z_2) = \| Z_1 - Z_2 \|_{C_\lambda}
\end{equation*}
\begin{comment}
It can be extended to the continuous case by
\begin{equation} \label{contGeneralizedCramer}
\begin{split}
	\ell_{\lambda}^2(Z_1,Z_2) = &\int_{x=z_1}^{z_K} \left( \int_{y=z_1}^{x} [ \Pi_{e^\perp} Z_1(y) - \Pi_{e^\perp} Z_2(y) ] dx \right)^2 dy \ + \\
	&\lambda \left( \int_{y=z_1}^{z_K} [Z_1(y) - Z_2(y)] dx \right)^2
\end{split}
\end{equation}
where for each return distribution $Z$
\begin{equation*}
	 \Pi_{e^\perp} Z(x) = Z(x) - \int_{y=z_1}^{z_K} Z(y) dy.
\end{equation*}
\end{comment}
Looking at Equation \ref{generalizedCramer}, the idea is that the first term on the right-hand side penalizes the difference in cumulative probabilities of both distributions, whereas the second term accounts for the difference in mass; we refer to \cite{DRLlinear} to get more details.

But, in addition to the Cramér distance generalization, authors in \cite{DRLlinear} also take into account the distribution $\xi$ according to which the states to be updated are sampled; %It can be whatever distribution, but in our case, given that we are following a policy $\pi$, we can consider it to be the stationary distribution $d_\pi$ (i.e. $d_\pi \in \Delta(\Sspace)$ such that $d_\pi = P_\pi d_\pi$; see a more detailed definition in Section \ref{COPTDsetting}).
they use it to define a \textit{$\xi$-weighted Cramér distance} over value distributions as
\begin{equation} \label{weightedCramer} 
	\ell_{\xi,\lambda}^2 (Z_1,Z_2) := \sum_{s \in \Sspace} \xi(s) \ell_{\lambda}^{2}  (Z_1(s),Z_2(s)).
\end{equation}
The two presented distances are precisely the key to all further results shown in \cite{DRLlinear}, and therefore to prove the desired convergence of CDRL algorithm with linear function approximation.

In particular, the previously commented projections of value distributions $\Pi_1$ and $\Pi_2$ are defined through $\ell_{\lambda}^2$ and $\ell_{\xi,\lambda}^2$, respectively. The former, which we will denote by $\Pi_{\lambda,\mathcal{P}}$, projects any value distribution $Z$ onto the subspace $\mathcal{P}$ according to
\begin{equation}
	\Pi_{\lambda,\mathcal{P}} Z := {\arg \min}_{Z'\in \mathcal{P}} \ell_{\lambda}^2 (Z, Z');
\end{equation} 
the previously considered Cramér projection $\Pi_C$ satifies this minimization expression\cite{DRLlinear}. The latter, which in our notation will be expressed as $\Pi_{\xi,\lambda, \Phi}$, performs a $\xi$-weighted projection of any value distribution $Z$ onto the set of value distributions $\Phi$ by:
\begin{equation}
	\Pi_{\xi,\lambda, \Phi} Z := {\arg \min}_{\Phi\Theta, \Theta \in \R^{m \times K}} \ell_{\xi,\lambda}^2 (Z, \Phi\Theta)
\end{equation} 
\begin{comment}
\begin{lemma}
	Let Z be an arbitrary value distribution supported on $\mathscr{P}([z_1,z_K])$. The $\xi$-weighted projection of $Z$ onto $\Phi$, $\Pi_{\xi,\lambda, \Phi} Z$, is equal to the $\xi$-weighted projection of $Z$ onto the parametric family $\mathcal{P}$, $\Pi_{\xi,\lambda, \mathcal{P}} Z$.
\end{lemma}
\end{comment}
The following Lemma \ref{projectionNonexpansion} shows an intewresting result regarding the good behaviour of this projection $\Pi_{\xi,\lambda, \Phi}$
\begin{lemma} \label{projectionNonexpansion}
	$\Pi_{\xi,\lambda, \Phi}$ is a non-expansion in $\ell_{\xi,\lambda}^2$; for every pair of return distributions $Z, Z'$, we have
	\begin{equation*}
		\ell_{\xi,\lambda}^2 (\Pi_{\xi,\lambda, \Phi} Z, \Pi_{\xi,\lambda, \Phi} Z') \leq \ell_{\xi,\lambda}^2(Z,Z')
	\end{equation*}
\end{lemma}

 At this point, before facing the proof of convergence of the whole CDRL algorithm, we would need to rewrite distances $\ell^2_\lambda$ and $\ell^2_{\xi,\lambda}$ into two separate components: along dimension $e_K$ and along the subspace $A$ orthogonal to $e_K$, $A:= \Pi_{e_K^\top}C$:
\begin{equation}
	\left\{
	\begin{array}{ll}
		\ell_{\lambda}^2(Z_1,Z_2) &= \| Z_1 - Z_2 \|^2_{AA^\top} + \lambda \| Z_1 - Z_2 \|^2_{e_K e_K^\top}  \\
		\ell_{\xi,\lambda}^2(Z_1,Z_2) &= \| Z_1 - Z_2 \|^2_{\xi,AA^\top} + \lambda \| Z_1 - Z_2 \|^2_{\xi, e_K e_K^\top} 
	\end{array}
	\right.
\end{equation}
The reason relies in Lemma \ref{linearBellmanContraction}, which states that the combined operator $\Pi_{\lambda,\mathcal{P}} \T^\pi$ behaves differently in each of these dimensions: while contracting all dimensions orthogonal to $e_K$ -i.e. the subspace $A$- by a factor $\gamma^{1/2}$, it is only a non-expansion along $e_K$. 

\begin{lemma} \label{linearBellmanContraction}
	Let $d_\pi$ the stationary distribution induced by policy $\pi$. For any two return distributions $Z_1,Z_2$, we have
	\begin{equation*}
		\begin{array}{rcl}
			\| \Pi_{\lambda,\mathcal{P}} \T^\pi Z_1 - \Pi_{\lambda,\mathcal{P}} \T^\pi Z_2 \|^2_{d_\pi, AA^\top} & \leq & \gamma \| Z_1 - Z_2 \|^2_{d_\pi, AA^\top} \\
			\| \Pi_{\lambda,\mathcal{P}} \T^\pi Z_1 - \Pi_{\lambda,\mathcal{P}} \T^\pi Z_2 \|^2_{d_\pi, e_KE_K^\top} & \leq & \| Z_1 - Z_2 \|^2_{d_\pi, e_KE_K^\top}
		\end{array}
	\end{equation*}
\end{lemma}

Note that, instead of the more general notation $\xi$, we are now specifying the stationary distribution $d_\pi$ (i.e. $d_\pi \in \Delta(\Sspace)$ such that $d_\pi = P_\pi d_\pi$; see a more detailed definition in Section \ref{offpolicySetting}). This is simply due to the fact that in our CDRL algorithm we are actually following policy $\pi$ to describe the agent trajectories we are learning from.

Finally, we present the theorem of \cite{DRLlinear} that provides us with the desired convergence guarantee for CDRL with Linear Function Approximation:


\begin{theorem} \label{convergenceLinearCDRL}
	Let $d_\pi$ the stationary distribution induced by policy $\pi$. The process
	\begin{equation*}
		Z_0 := \Phi \Theta_0, \quad Z_{k+1} := \hat{\Pi}_{d_\pi,\lambda,\Phi} T^\pi Z_k
	\end{equation*}
	converges to a set $S$ such that, for any two $Z, Z' \in S$, there is a $\Sspace$-indexed vector of constants $\alpha$ such that
	\begin{equation*}
		Z(s) = Z'(s) + \alpha(s)e_K.
	\end{equation*}
	If $\lambda > 0$, $S$ consists of a single point $\hat{Z}$ wich is the fixed point of the process. Moreover, we can bound the error of this fixed point with respect to the true return distribution $Z^\pi$ by
	\begin{equation*}
		\ell^2_{d_\pi, \lambda} (\hat{Z}, Z^\pi) \leq \frac{1}{1-\gamma} \ell^2_{d_\pi, \lambda} (\Pi_{d_\pi,\lambda,\Phi} Z^\pi, Z^\pi) - \frac{\gamma \lambda}{1-\gamma} \| \hat{Z} - Z^\pi \|_{d_\pi, e_Ke_K^\top},
	\end{equation*}
	where the second terms measures the difference in mass between $\hat{Z}$ and $Z^\pi$.
\end{theorem}

OPERATOR $\hat{\Pi}_{d_\pi,\lambda,\Phi}$, BASED ON LOSS
\begin{equation} \label{generalizedCramerLoss}
	\hat{\ell}_{\lambda}^2(Z_1,Z_2) = (Z_1 - Z_2)^\top \Pi_{e_K^\top} C C^\top \Pi_{e_K^\top} (Z_1 - Z_2) + \lambda (Z_2 e -1)^2 
\end{equation}
with the explicit normalization penalty $(Z_2 e -1)^2$ that encourages return distributions to have unit mass. Not a distance!! $\xi$-weighted Cramér loss $\hat{\ell}_{\xi, \lambda}^2$ is accordingly derived, and operator $\hat{\Pi}_{d_\pi,\lambda,\Phi}$, for a return distribution $Z$, finds the value distribution in the span of $\Phi$ which minimizes $\hat{\ell}_{\xi, \lambda}^2(Z, \cdot)$:
\begin{equation} \label{operatorCramérLoss}
	\hat{\Pi}_{d_\pi,\lambda,\Phi} Z = \Phi \Theta^* \ \textnormal{ where } \ \Theta^* = {\arg \min}_{\Theta} \hat{\ell}_{\xi, \lambda}^2(Z, \Phi \Theta)
\end{equation}

As pointed out in \cite{DRLlinear}, the parameter $\lambda$ plays an important role in the theorem, both to guarantee convergence and to bound the approximation error. At a high level, this makes sense: a high value of $\lambda$ forces the algorithm to output something close to a distribution, at the expense of actual predictions. On the other hand, taking $\lambda=0$ yields a process which may not converge to a single point.
 
 \setlength{\algomargin}{1.5em}
\SetAlCapHSkip{0em}
\begin{algorithm}[H]
\caption{CDRL with Linear Function Approximation} \label{linearCDRL}
\DontPrintSemicolon
\SetKwComment{CustomComment}{\#}{}
\SetKwInput{KwIn}{\hspace{-1.5em} Input}
\SetKwInput{KwOut}{\hspace{-1.5em} Output}
\SetKwInOut{Require}{\hspace{-1.5em} Require}
\Require{$Z_{\Theta_t}(s) = \Theta^\top_t \phi(s)$ for each $s\in \Sspace$}
\KwIn{A state $s$}
	\CustomComment{Compute distributional Bellman target}
    \uIf{\textnormal{Categorical Policy Evaluation}}{
    	$\widehat{Z}(s) \leftarrow \T^\pi Z_{\Theta_t}(s)$\;}
    \ElseIf{\textnormal{Categorical Q-learning}}{
    	$\widehat{Z}(s) \leftarrow \T^* Z_{\Theta_t}(s)$\;}
    \CustomComment{Project target onto support}
    $\widehat{Z}_{*}(s) \leftarrow \Pi_C \widehat{Z}(s)$\;
    \CustomComment{Project target onto space $\mathcal{L}$}
    $\widehat{Z}_{t}(s) \leftarrow \Pi_\Phi \widehat{Z}_{*}(s)$\;
    \CustomComment{Update weight matrix}
    $\Theta_{t+1} \leftarrow f(\Theta_t, \widehat{Z}_{t}(s) - Z_{\Theta_t}(s))$ \;
    \textnormal{Generate new estimate }$Z_{\Theta_{t+1}}(s) = \Theta^\top_{t+1} \phi(s)$ \;
\KwOut{Estimate $Z_{\Theta_{t+1}}(s)$ for each $s \in \Sspace$}

\end{algorithm}

\newpage


\thispagestyle{plain}

\thispagestyle{plain}
\section{Covariate Shift Ratio}

TODO: INTRO off-policy learning and COPTD

Original paper \cite{COPTD}:

The problem of on-line off-policy evaluation (OPE) has been actively studied in the last decade due to its importance both as a stand-alone problem and as a module in a policy improvement scheme. However, most Temporal Difference (TD) based solutions ignore the discrepancy between the stationary distribution of the behavior and target policies and its effect on the convergence limit when function approximation is applied. In this paper we propose the Consistent Off-Policy Temporal Difference (COP-TD(lambda, beta)) algorithm that addresses this issue and reduces this bias at some computational expense. We show that COP-TD(lambda, beta) can be designed to converge to the same value that would have been obtained by using on-policy TD(lambda) with the target policy. Subsequently, the proposed scheme leads to a related and promising heuristic we call log COP-TD(lambda, beta). Both algorithms have favorable empirical results to the current state of the art online OPE algorithms. Finally, our formulation sheds some new light on the recently proposed Emphatic TD learning.

Reinforcement Learning (RL) techniques were successfully applied in fields such as robotics, games, marketing and more (Kober et al., 2013; Al-Rawi et al., 2015; Barrett et al., 2013). We consider the problem of offpolicy evaluation (OPE) – assessing the performance of a complex strategy without applying it. An OPE formulation is often considered in domains with limited sampling capability. For example, marketing and recommender systems (Theocharous  Hallak, 2013; Theocharous et al., 2015) directly relate policies to revenue. A more extreme example is drug administration, as there are only few patients in the testing population, and sub-optimal policies can have life threatening effects (Hochberg et al., 2016). OPE can also be useful as a module for policy optimization in a policy improvement scheme (Thomas et al., 2015a).In this paper, we consider the OPE problem in an on-line setup where each new sample is immediately used to update our current value estimate of some previously unseen policy. We propose and analyze a new algorithm called COP-TD.

Our algorithm resembles \cite{ETD} Emphatic TD that was extended by \cite{generalizedETD} to the general parametric form ETD(lambda, beta). We clarify the connection between the algorithms and compare them empirically. Unlike ETD(lambda, beta), COP-TD(lambda,beta)’s effectiveness depends on the available resources. The number of features can be adjusted accordingly to provide the most affordable approximation. The added cost is fine-tuning another step- size, though beta’s effect is less prominent.


Conclusions: Research on off-policy evaluation has flourished in the last decade. While a plethora of algorithms were suggested so far, ETD(lambda, beta) by Hallak et al. (2015) has perhaps the simplest formulation and theoretical properties. Unfortunately, ETD(lambda, beta) does not converge to the same point achieved by on-line TD when linear function approximation is applied.We address this issue with COP-TD(lambda,beta) and proved it can achieve consistency when used with a correct set of features, or at least allow trading-off some of the bias by adding or removing features. Despite requiring a new set of features and calibrating an additional update function, COP-TD(lambda,beta)’s performance does not depend as much on beta as ETD( lambda, beta), and shows promising empirical results.

Paper Carles Discounted COP-TD \cite{DCOPTD}:

Central to reinforcement learning is the idea that an agent should learn from experience. While many algorithms learn in a purely online fashion, sample-efficient methods typically make use of past data, viewed either as a fixed dataset, or stored in a replay memory (Lin 1993, Mnih et al. 2015). Because this past data may not be generated according to the policy currently under evaluation, the agent is said to be learning off-policy (Sutton and Barto 2018).

By now it is well-documented that off-policy learning may carry a signifcant cost when combined to function approximation. Early results have shown that estimating the value function off-policy, using Bellman updates, may diverge (Baird 1995), (Tsitsiklis and Van Roy 1997). More recently, value divergence was perhaps the most significant issue dealt with in the design of the DQN agent (Mnih et al. 2015), and remains a source of concern in deep reinforcement learning (van Hasselt, Guez, and Silver 2016).

Further, under off-policy learning, the quality of the Bellman fixed point suffers as studied by Kolter (2011) and Munos (2003). The value function error can be unboundedly large even if the value function can be perfectly approximated. Hence, even in the case where convergence to the fixed point with off-policy data occurs, solutions can be of poor quality. Thus, the existing TD learning algorithms with convergence guarantees under off-policy data (Maei et al. 2009), (Sutton et al. 2009) can still suffer from off-policy issues.

This paper studies the covariate shift method for dealing with the off-policy problem. The covariate shift method, studied by Hallak and Mannor (2017) and Sutton, Mahmood, and White (2016), reweights online updates according to the ratio of the target and behavior stationary distributions. Under optimal conditions, the covariate shift method recovers convergent behavior with linear approximation, breaking what Sutton and Barto (2018) call the “deadly triad” of reinforcement learning. We argue the method is particularly appealing in the context of replay memories, where the reweighting can be replaced by a reprioritization scheme similar to that of Schaul et al. (2016).

We improve on Hallak and Mannor’s COP-TD algorithm, which has provable guarantees but is dificult to implement in a deep reinforcement learning setting. First, we introduce a discount factor into their update rule to obtain a more stable algorithm. Second, we develop an alternative normalization scheme that can be combined with deep networks, avoiding the projection step necessary in the original algorithm.

Conclusion: In this paper we revisited Hallak and Mannor’s COP-TD algorithm and extended it to be applicable to the deep reinforcement learning setting. While these results on the Atari 2600 suite of games remain preliminary, they demonstrate the practicality of learning the covariate shift in complex settings. We believe our results further open the door to increased sample efficiency in deep reinforcement learning.

We emphasize that the instabilities observed when learning the covariate shift under prioritized sampling point to the importance of the data distribution used to learn the ratios. Which distribution is optimal will be the focus of future work. The covariate shift method is a “backward” off-policy method, in the sense that it corrects a mismatch between distributions based on past transitions. It would be interesting to combine our method to “forward” off-policy methods such as Retrace (Munos et al. 2016), which have also yielded good results on the Atari 2600 (Gruslys et al. 2018). Then, it would be interesting to understand whether overfitting does occur due to a smaller effective replay size, and how this can be addressed. Finally, an exciting avenue would be extending the method to the more general case where multiple policies have generated off-policy data, which would allow COP-TD to be applied in the standard control setting.


\subsection{Off-Policy Learning Setting} \label{offpolicySetting}

As stated in \cite{DCOPTD}, here we move to the \textit{policy evaluation} problem within \textit{off-policy learning}, where we want to learn the value function $V^\pi$ of a \textit{target policy} $\pi$ from samples drawn from $P$ and a \textit{behaviour policy} $\mu$. 

We first introduce some useful notation:
\begin{itemize}
    \item The Bellman equation for the state value function can be expressed in vector notation as $V^\pi = r_\pi + \gamma P_\pi V^\pi$, where $V^\pi \in \R^{n}$, $r_\pi \in \R^n$ and $P_\pi \in \R^{n \times n}$. The value function is in fact the fixed point of the \textit{Bellman operator} $\mathcal{T}_\pi: \R^{n \times n} \rightarrow \R^{n \times n}$, defined as $\mathcal{T}_\pi V := r_\pi + \gamma P_\pi V $. It defines a single step of \textit{bootstrapping}: the process $V^{k+1} := \mathcal{T}_\pi V^k$ converges to $V^\pi$.
    \item Let $d \in \R^n$; we write $D_d \in \R^{n\times n}$ for the corresponding diagonal matrix, and consider the weighted squared seminorm notation of vectors $x\in\R^n$ $||x||^2_{A} := ||A x||^2 = x^T A^T A x$, $||x||^2_d :=||x||^2_{D_d}= \sum_{i=1}^n d(i)^2 x(i)^2$.
    \item $e\in\R^n$ accounts for the vector of all ones, and $\Delta(\Sspace)$ for the simplex over states: $d \in \Delta(\Sspace) \implies d^T e = 1, d \geq 0$.
    \item $d \in \Delta(\Sspace)$ is the stationary distribution of a transition function $P$ if and only if $d=d \cdot P$. This distribution is unique when $P$ defines a Markov chain with a single recurrent class\cite{unichain}.
\end{itemize}
In this particular setting we distinguish between two different state-to-state transition functions, $P_\pi$ and $P_\mu$, one for each policy; their respective stationary distributions will be represented by $d_\pi$ and $d_\mu$. 

\subsubsection*{Linear Function Approximation}
In most RL problems, it is unfeasible to learn the value of each state separately due to the large state spaces involved; in these cases, we get value function estimates through \textit{function approximators}. Among them, it is common to consider a \textit{linear function approximation}, which uses a mapping from states to features $\phi: \Sspace \rightarrow \R^k$. In these cases, the approximate value function at a state $s$ can be expressed as the inner product of a feature vector with a vector of weights $\theta \in \R^k$:
\begin{equation} \label{linearApproxSV}
    \Hat{V}(s) = \phi(s)^T \theta
\end{equation}
which can be written as $\Hat{V} = \Phi \theta$ in vector notation, being $\Phi \in \R^{n\times k}$ the matrix of row vectors. The \textit{semi-gradient update rule} for TD learning\cite{TD} learns an estimation of $V^\pi$ from sample transitions. Given a starting state $s\in\Sspace$, a successor state $s'\sim P_\pi(\cdot | s)$, and a step-size parameter $\alpha > 0$, this update is
\begin{equation} \label{updateTD}
    \theta \leftarrow \theta + \alpha \left[ r_\pi(s) + \gamma\phi(s')^T \theta - \phi(s)^T \theta \right] \phi(s)
\end{equation}

The expected behaviour of this update rule is described by the \textit{projected Bellman operator} $\Pi_d \mathcal{T}_\pi$, a combination of the usual Bellman operator with a projection $\Pi_d$ in norm $||\cdot||_d$ -for some $d\in\Delta(\Sspace)$- onto the span of $\Phi$\cite{projectedBellman}. In fact, the stationary point of \ref{updateTD}, if it exists, is the solution of the \textit{projected Bellman equation} $\Hat{V}^\pi = \Pi_d \mathcal{T}_\pi \Hat{V}^\pi$. If the projection is performed under $d_\pi$, we can guarantee convergence to a fixed point\cite{rlformulation}; otherwise -i.e. $d \neq d_\pi$- it may lead to divergence.

We want to note that the proof of convergence when $d=d_\pi$ is done by showing that the combined operator $\Pi_{d_\pi} T^\pi$ is a contraction in the $d_\pi$-weighted norm, i.e. for any two value functions $V,V' \in \R^n$
\begin{equation*}
	\| \Pi_{d_\pi} T^\pi V - \Pi_{d_\pi} T^\pi V' \|_{d_\pi} \leq \gamma \| V-V' \|_{d_\pi}
\end{equation*}
From this, it only remains to apply Banach's fixed point theorem. 
Furthermore, as presented in \cite{DCOPTD}, not only $d=d_\pi$ provides convergence, but this choice also seems optimal in terms of the quality of the fixed point under off-policy data with linear function approximation.


\subsection{Covariate Shift Approach} \label{COPTD}

Supposing that stationary distributions $d_\pi$ and $d_\mu$ are known, and that states are updated according to $s\sim d_\mu$, the covariate shift approach presented in \cite{COPTD} uses importance sampling to redefine \ref{updateTD} so that the semi-gradient update rule can be considered \textit{under the sampling distribution} $d_\pi$:
\begin{equation} \label{updateTDimportance}
    \theta \leftarrow \theta + \alpha \frac{d_\pi(s)}{d_\mu (s)}\left[ r(s,a) + \gamma\phi(s')^T \theta - \phi(s)^T \theta \right] \phi(s)
\end{equation}
with $a\sim \mu(\cdot|s)$, $s' \sim P(\cdot|s,a)$ as before. 

Hence, the Consistent Off-Policy Temporal Difference\cite{COPTD} (COP-TD) algorithm seek to learn that covariate shift ratio $d_\pi / d_\mu$ from samples by bootstrapping from a previous prediction, similar to temporal difference learning. Given a step size $\alpha > 0$, a ratio vector $c \in \R^n$ and a sample transition $(s_t,a_t,s_{t+1}) = (s,a,s')$ drawn from $d_\mu$, $\mu(\cdot|s)$ and $P(\cdot|s,a)$, respectively, the COP-TD update is
\begin{equation} \label{COP}
    c(s') \leftarrow c(s') + \alpha \left[\frac{\pi(a|s)}{\mu(a|s)}c(s) + c(s') \right]
\end{equation}
The expected behaviour of this learning rule, which learns "in reverse" compared to TD learning, is captured by the \textit{COP operator} $Y$:
\begin{equation} \label{COPoperator}
    (Yc)(s') := \E_{s \sim d_\mu, a \sim \mu(\cdot|s)} \left[ \frac{\pi(a|s)}{\mu(a|s)}c(s)  \bigg| s' \right]
\end{equation}

Note that the condition $s_{t+1} = s'$ in the expectation of \ref{COPoperator} forces to take into account the distribution of previous state-action pairs $(s,a)$ according to policy $\mu$. The distribution of the possible previous states $s$ is given by the time-reversal transition function $\Bar{P}_\mu$, whose entries are:
\begin{equation}
\begin{split}
    \Bar{P}_\mu (s | s') &:= Prob_\mu (s_t = s | s_{t+1} = s') \\
    &= \frac{Prob_\mu ( s_{t+1} = s' | s_t = s) Prob_\mu(s_t=s)}{Prob_\mu ( s_{t+1} = s')} \\
    &= \frac{P_\mu (s' | s) d_\mu(s)}{d_\mu(s')} 
\end{split}
\end{equation}
Or, equivalently, $ \Bar{P}_\mu = D_{d_\mu}^{-1}P_\mu^T D_{d_\mu}$ in vector notation. Regarding the distribution of the possible actions that lead to $s'$ from a certain state $s$ by following policy $\mu$, it will be represented by function $\Bar{\mu}$:  
\begin{equation}
\begin{split}
        \Bar{\mu}(a|s,s') &:= Prob_\mu(a_t=a | s_t = s, s_{t+1}=s') \\
        &= \frac{Prob_\mu (a_t=a, s_t = s, s_{t+1}=s')}{Prob_\mu (s_t = s, s_{t+1}=s')} \\
        %&= \frac{Prob_\mu (s_{t+1}=s' | a_t=a, s_t = s) Prob_\mu(a_t=a, s_t = s)}{Prob_\mu (s_t = s, s_{t+1}=s')} \\
        &= \frac{Prob_\mu (s_{t+1}=s' | a_t=a, s_t = s) Prob_\mu (a_t=a | s_t = s) Prob_\mu (s_t = s)}{Prob_\mu (s_{t+1}=s' | s_t = s) Prob_\mu (s_t = s)} \\
        &= \frac{P(s'|s,a)\mu(s|a)}{P_\mu (s'|s)}
\end{split}
\end{equation}


Bearing in mind the introduced notation, the expectation in \ref{COPoperator} can be rewritten and expanded:
\begin{equation} \label{expansionCSR}
\begin{split}
    \E_{s \sim d_\mu, a \sim \mu(\cdot|s)} \left[ \frac{\pi(a|s)}{\mu(a|s)}c(s)  \bigg| s' \right] &= \E_{s \sim \Bar{P}_\mu (\cdot | s'), a \sim \Bar{\mu}(\cdot|s,s')} \left[ \frac{\pi(a|s)}{\mu(a|s)}c(s) \right] \\
    &= \sum_s \Bar{P}_\mu (s | s') \sum_a \Bar{\mu}(a|s,s') \frac{\pi(a|s)}{\mu(a|s)} c(s) \\
    &= \sum_s \left( \frac{P_\mu (s' | s) d_\mu(s)}{d_\mu(s')}  \right) \sum_a  \left( \frac{P(s'|s,a)\mu(s|a)}{P_\mu (s'|s)} \right) \frac{\pi(a|s)}{\mu(a|s)} c(s) \\
    &= \frac{1}{d_\mu(s')} \sum_{s} d_\mu(s) c(s) \sum_a \pi(a|s) P(s'|s,a) \\
    &= \frac{1}{d_\mu(s')} \sum_{s} P_\pi(s'|s) d_\mu(s) c(s) \\
    %&=  \frac{1}{d_\mu(s')} \E_{s\sim d_\mu} \left[P_\pi(s'|s) c(s)\right]
\end{split}
\end{equation}
Thus, according to the last term of \ref{expansionCSR}, the COP operator $Y$ can be expressed in vector notation as
\begin{equation}
    Yc = D_{d_\mu}^{-1}P_\pi^T D_{d_\mu} c
\end{equation}
Regarding the behaviour of the operator, the following result of \cite{DCOPTD} guarantees that the process $c^{k+1} := Yc^k$ converges, and that any multiple of $d_\pi / d_\mu$ is a fixed point:
\begin{theorem} \label{convergenceCOPTD}
	Suppose that $P_\pi$ defines an ergodic Markov chain on $\Sspace$, and let $c^0 \in \Delta(\Sspace)$. Then the process $c^{k+1} := Yc^k$ converges to $C \frac{d_\pi}{d_\mu}$, $C \in \R^+$
\end{theorem}

Furthermore, we can reduce the set of fixed points to the covariate shift ratio exclusively by considering a normalized version of the COP operator:
\begin{corol}
	Suppose conditions of Theorem \ref{convergenceCOPTD} are satisfied. Let
	\begin{equation*}
		(\bar{Y} c)(s') := \frac{(Yc)(s')}{\sum_{s} (Yc)(s)}
	\end{equation*}
	be the normalized COP operator. Then the unique fixed point of $\bar{Y}$ is the ratio $d_{\pi}/d_\mu$, to which $c^{k+1} = \bar{Y} c^k$ converges.
\end{corol}

\subsubsection*{COP-TD with Linear Function Approximation}

Whenever the value function is approximated, we should expect to learn a ratio $\hat{c}$ estimate as well. In our analysis, we consider a linear approximation of the form
\begin{equation*}
	\hat{c}(s) = \phi(s)^\top w
\end{equation*}
where again $\phi$ defines a map from states to features, $\phi: \Sspace \rightarrow \R^k$, and now $w \in \R^k$ is the vector of weights we are interested in learning. Note that in such a model negative ratio values are allowed; this can be avoid in practice by clipping those values at zero.

The introduced linear model, given a sample transition $(s,a,s')$ drawn from $d_\mu$, $\mu(\cdot|s)$ and $P(\cdot|s,a)$, respectively, induces the following semi-gradient update
\begin{equation} \label{updateLinearCOPTD}
	\tilde{w} \leftarrow w + \alpha \left[ \frac{\pi(a|s)}{\mu(a|s)} \phi(s)^\top w - \phi(s') w \right] \phi(s'),
\end{equation}
%which is similar to \ref{updateTD} but implementing \ref{COP}. 
which can be thought as a $d$-weighted projection $\Pi_d$ for some $d \in \Delta(\Sspace)$\cite{DCOPTD}. However, an additional step is required for granting that the resulting ratio estimate $\hat{c}$ corresponds to some proper distribution ratio $d / d_{\mu}$ for that $d \in \Delta(\Sspace)$, just as we want to. This is solved in \cite{COPTD} by following the update rule by a projection onto the $d_\mu$-weighted simplex $\Delta_{\Phi,d_\mu}$
\begin{equation}
	w \leftarrow {\arg \min}_{u\in W_{\Phi,d_\mu}} \| u - \tilde{w} \|
\end{equation}
where $ W_{\Phi,d_\mu} := \{ u \in \R^k : \sum_{s \in \Sspace} d_\mu(s) \phi(s)^\top u = 1, \phi(s)^\top u \geq 0 \}$. Looking at the definition of $W_{\Phi,d_\mu}$, we can identify this second projection $\Pi_{\Delta_{\Phi,d_\mu}}$ as a normalization step as well.

The following Lemma \ref{normalizationLinearCOPTD} of \cite{DCOPTD} shows that, in fact, the normalization component of operator $\Pi_{\Delta_{\Phi,d_\mu}}$ is not only a convenience but also a requisite to get a good convergence guarantee. 
\begin{lemma} \label{normalizationLinearCOPTD}
	Let $Y$ be a symmetric COP-TD operator and $\Pi$ the projection onto $\Phi$ in $L_2$ norm. If $d_\pi / d_\mu$ is not in the span of $\Phi$, then $c=0$ is the only fixed point of the process $c^{k+1} = \Pi Y c^k$.
\end{lemma}
\noindent Hence, in order to get a meaningful approximation ratio, we must consider the repeated application of the combined operator $\Pi_{\Delta_{\Phi,d_\mu}} \Pi_d Y$, so that iterating $\hat{c}^{k+1} := \Pi_{\Delta_{\Phi,d_\mu}} \Pi_d Y \hat{c}^k$ provides the estimate.

In practice, however, that combination of the COP operator with the projection onto the $d_\mu$-weighted simplex can be really hard to implement; authors in \cite{COPTD} presented a method for approximating the projection step in an online, sample-based setting for linear function approximation, but to the best of our knowledge there are no analogous implementation designs for other kind of function approximations, like neural networks.





\subsection{Discounted COP-TD}

Despite the good properties shown of COP-TD, there are two practical limitations of the algorithm that specially motivated the authors in \cite{DCOPTD} to work in an improvement. On the one hand, we have the previously commented difficulties to practically implement the projection operator $\Pi_{\Delta_{\Phi,d_\mu}}$ with function approximations other than linear ones. On the other hand, we note that the COP operator $Y$ lacks a result ensuring it is a contraction mapping, which can make the algorithm converge at a slow rate or with high variance, or even be unstable together with function approximation.

They address these two issues in \cite{DCOPTD} through the $\hat{\gamma}$\textit{-discounted COP-TD} learning rule; given a step size $\alpha > 0$, a ratio vector $c \in \R^n$ and a sample transition $(s_t,a_t,s_{t+1}) = (s,a,s')$ drawn from $d_\mu$, $\mu(\cdot|s)$ and $P(\cdot|s,a)$, respectively, it modifies the previous COP-TD update \ref{COP} by
\begin{equation} \label{DCOP}
    c(s') \leftarrow c(s') + \alpha \left[\hat{\gamma} \frac{\pi(a|s)}{\mu(a|s)}c(s) + (1-\hat{\gamma}) - c(s') \right]
\end{equation}
Again, the expected behaviour of this rule is captured by the corresponding $\hat{\gamma}$-discounted COP operator $Y_{\hat{\gamma}}$
\begin{equation} \label{DCOPoperator}
    (Y_{\Hat{\gamma}}c)(s') := \E_{s \sim d_\mu, a \sim \mu(\cdot|s)} \left[ \hat{\gamma} \frac{\pi(a|s)}{\mu(a|s)}c(s) + (1-\hat{\gamma}) \bigg| s' \right] = \hat{\gamma} (Yc)(s') + (1-\hat{\gamma}),
\end{equation}
where $Y$ is the COP operator (\ref{COPoperator}). In vector notation we have
\begin{equation}
    Y_{\Hat{\gamma}}c = \Hat{\gamma} D_{d_\mu}^{-1}P_\pi^T D_{d_\mu} c + (1-\Hat{\gamma})e
\end{equation}
From the above expressions, it is straightforward to notice that we recover COP-TD for $\hat{\gamma} = 1$, i.e. $Y_1=Y$.










\newpage

\thispagestyle{plain}
\section{Distributional COP-TD}

TODO: INTRO

Now we are interested in going beyond the notion of value and consider the estimation of the CS ratio from a distributional perspective, similarly to what was done in \cite{DRL} within the reinforcement learning setting. Our starting point could be 
\begin{equation} \label{distCSR}
    X(s') \stackrel{D}{=} \frac{\pi(A^\mu_{s,s'}|S^\mu_{s'})}{\mu(A^\mu_{s,s'}|S^\mu_{s'})} X(S^\mu_{s'}) 
    %X(s') \stackrel{D}{=} \hat{\gamma} \frac{\pi(A^\mu_{s,s'}|S^\mu_{s'})}{\mu(A^\mu_{s,s'}|S^\mu_{s'})} X(S^\mu_{s'}) +1 -\hat{\gamma}
\end{equation}
where
\begin{itemize}
    \item $X$ is the random ratio between distributions of achieving a certain state, its expectation being the covariate shift ratio 
    \begin{equation}
        \frac{d_\pi}{d_\mu}(s) = c(s) = \E [X(s)]
    \end{equation}
    
    \item $(S^\mu_{s'},A^\mu_{s,s'} )$ is the previous state-action random variable:
   	\begin{itemize}
    	\item The random variable $S^\mu_{s'}$ represents the states from which state $s'$ is achievable by following policy $\mu$; $S^\mu_{s'} =s$ with probability $\Bar{P}_\mu(s|s')$
   	 	\item $A^\mu_{s,s'}$ encodes the random action that can be taken to get state $s'$ from a state $s\sim S^\mu_{s'}$ according to policy $\mu$, so $A^\mu_{s,s'} =a$ with probability $\Bar{\mu}(a|S^\mu_{s'},s')$
   	 \end{itemize}
 
\end{itemize}

Recalling equation \ref{distCSR}, note that it intrinsically expresses the random ratio of a state $s_{t+1}$, $X(s_{t+1})$, as a mixture distribution with the 'corrected' previous state random ratios $\rho(s_{t},a_{t}) X(s_{t})$ as mixing components, and the previous state-action random variable $(S^\mu_{s_{t+1}},A^\mu_{s_{t},s_{t+1}} )$ as the mixing distribution:
\begin{equation} \label{densityDCS}
    f_{X(s_{t+1})}(x) = \sum_{s_{t}} \Bar{P}_\mu (s_{t} | s_{t+1}) \sum_{a_t} \Bar{\mu}(a_t|s_{t},s_{t+1}) f_{\frac{\pi(s_{t},a_{t})}{\mu(s_{t},a_{t})} X(s_{t})}(x) 
\end{equation}

\begin{notation}
	So as to reduce the complexity and increase the readability of the formulation, we introduce the following notation:
	
	\begin{itemize}
		\item Let define $\rho$ the policy ratio $\pi / \mu$:
%\begin{equation}
%	\rho(s_{t},a_{t}) = \rho (A^\mu_{s_{t},s_{t+1}},S^\mu_{s_{t+1}}) := \frac{\pi(A^\mu_{s_{t},s_{t+1}}|S^\mu_{s_{t+1}})}{\mu(A^\mu_{s_{t},s_{t+1}}|S^\mu_{s_{t+1}})}
%\end{equation}
		\begin{equation*}
			\rho(a,s) := \frac{\pi(a|s)}{\mu(a|s)}
		\end{equation*}

	\item Note in equation \ref{densityDCS} that there are as many mixture components as state-action pairs $(s_t,a_t)$; thus, we can iterate the summation over all these possible pairs and define each corresponding mixture weight $\alpha$ as
		\begin{equation*}
			\alpha(s_{t},a_{t}; s_{t+1}) = \Bar{P}_\mu (s_{t} | s_{t+1}) \Bar{\mu}(a_t|s_{t},s_{t+1}) 
		\end{equation*}
	\end{itemize}
\end{notation}

Considering the previous notation, the expression of the density function \ref{densityDCS} can be rewritten and expanded in the following way:
\begin{equation}
\begin{split}
    f_{X(s_{t+1})}(x) &= \sum_{(s_{t},a_{t})} \alpha(s_{t},a_{t};s_{t+1}) \cdot f_{\rho(s_{t},a_{t}) X(s_{t})}(x) \\
    &= \sum_{(s_{t},a_{t})} \alpha(s_{t},a_{t};s_{t+1}) \cdot f_{\rho(s_{t},a_{t}) \rho(S^{\mu}_{s_{t}},A^{\mu}_{s_{t-1},s_{t}}) X(S^{\mu}_{s_t}) )}(x) \\
    &= \sum_{(s_{t},a_{t})} \alpha(s_{t},a_{t};s_{t+1}) \cdot \sum_{(s_{t-1},a_{t-1})} \alpha(s_{t-1},a_{t-1};s_{t})   \cdot f_{\rho(s_{t},a_{t}) \rho(s_{t-1},a_{t-1}) X(s_{t-1})}(x) \\
    &= \sum_{(s_{t},a_{t})} \alpha(s_{t},a_{t};s_{t+1}) \cdots \sum_{(s_{0},a_{0})} \alpha(s_{0},a_{0};s_{1}) \cdot f_{ \left(\prod_{i=t}^{0} \rho(s_{i},a_{i})\right)  X(s_{0})}(x)
\end{split}
\end{equation}

Regarding its expectation, we can see that:
\begin{equation}
\begin{split}
    \E[X (s_{t+1})] &= \int_{-\infty}^{\infty} x f_{X(s_{t+1})}(x) dx \\
    &= \sum_{(s_{t},a_{t})} \alpha(s_{t},a_{t};s_{t+1}) \cdot \int_{-\infty}^{\infty} x f_{\rho(s_{t},a_{t}) X(s_{t})}(x) dx \\
    &= \sum_{(s_{t},a_{t})} \alpha(s_{t},a_{t};s_{t+1}) \cdot \E \left[ \rho(s_{t},a_{t}) X(s_{t}) \right] \\
    &=  \sum_{(s_{t},a_{t})} \alpha(s_{t},a_{t};s_{t+1}) \cdot \rho(s_{t},a_{t}) \cdot \E \left[X(s_{t}) \right] 
\end{split}
\end{equation}

\subsubsection*{Discounted Distributional CS Ratio}
Having already develop the previous formulation, it is straightforward to turn equation \ref{distCSR} into its discounted version as it is done in \cite{DCOPTD} in the value-based case:
\begin{equation} \label{distDCSR}
    X(s_{t+1}) \stackrel{D}{=} \hat{\gamma} \ \rho(S^{\mu}_{s_{t+1}},A^{\mu}_{s_{t},s_{t+1}}) X(S^{\mu}_{s_{t+1}}) + 1 -\hat{\gamma}
\end{equation}
as well as compute its corresponding expectation:
\begin{equation}
\E[X (s_{t+1})] = 1- \hat{\gamma}\left(1 - \sum_{(s_{t},a_{t})} \alpha(s_{t},a_{t};s_{t+1}) \cdot \rho(s_{t},a_{t}) \cdot \E \left[X(s_{t}) \right]  \right)
\end{equation}

Given that $\E \left[X(s) \right] = c(s)$, the previous equation is equivalent to the result of applying of the DCOP operator, whose convergence is analyzed and proved in \cite{DCOPTD}. This can help us prove the convergence to a fixed point in the distributional setting, whose expectation matches with the covariate shift ratio estimate $c$ in the value-based case.


\newpage

\thispagestyle{plain}
\section{Logarithmic Approach to Distributional COP-TD}

Note that the Distributional Covariate Shift equation is purely multiplicative, and so it is the associated update rule in the learning setting. 


Let's consider
\begin{equation}
	Y(s_{t+1}) := \log \left( X(s_{t+1}) \right) = \log \left( \rho(S^{\mu}_{s_{t+1}},A^{\mu}_{s_{t},s_{t+1}}) \right) + Y(S^{\mu}_{s_{t+1}}) 
\end{equation}
Its density function being
\begin{equation}
\begin{split}
	f_{Y(s_{t+1})} (x) &= \sum_{(s_{t},a_{t})} \alpha(s_{t},a_{t};s_{t+1}) \cdot f_{\log\left( \rho(s_{t},a_{t}) X(s_{t}) \right)}(x) \\
	&= \sum_{(s_{t},a_{t})} \alpha(s_{t},a_{t};s_{t+1}) \cdot f_{\log\left( \rho(s_{t},a_{t})\right) + Y(s_t)}(x)
\end{split}
\end{equation}

Recalling \textbf{Property \ref{mixtureProperty3}} of Mixture Distributions, we can express the expectation of $Y$ as:
\begin{equation}
\begin{split}
	\E[Y(s_{t+1})] &= \sum_{(s_{t},a_{t})} \alpha(s_{t},a_{t};s_{t+1}) \cdot \E \left[ \log \left( \rho(s_{t},a_{t}) X(s_{t}) \right) \right] \\
	&= \sum_{(s_{t},a_{t})} \alpha(s_{t},a_{t};s_{t+1}) \cdot \E \left[ \log \left( \rho(s_{t},a_{t}) \right)  + Y(s_t) \right] \\
	&= \sum_{(s_{t},a_{t})} \alpha(s_{t},a_{t};s_{t+1}) \cdot \left( \log \left( \rho(s_{t},a_{t}) \right) + \E [Y(s_t)] \right)
\end{split}
\end{equation}

TODO: ADDITIVE FIXED POINT
\\

And if
\begin{equation}
	U(s') = \exp (Y(s'))
\end{equation}
expectation:
\begin{equation}
\begin{split}
	\E[U(s_{t+1})] &= \int_{-\infty}^{\infty} \exp(y) f_{Y(s_{t+1})} (y) dy \\
	&= \int_{-\infty}^{\infty} \exp(\log(x)) f_{X(s_{t+1})} (x) dx \\
	&=  \int_{-\infty}^{\infty} x f_{X(s_{t+1})} (x) dx  \\
	&= \E [X(s_{t+1})]
\end{split}
\end{equation}


\newpage

\thispagestyle{plain}
\section{Implementation}

\subsection{Categorical COP-TD}

Analogously to \cite{DRL}, we can attempt to model the CS ratio distribution through a discrete parametric distribution with a certain set of atoms $\{x_i\}_{0 \leq i<M}$, $M \in \mathbb{N}$, as its support. The atom probabilities would be given by a parametric model $\theta : \mathcal{X} \rightarrow \mathbb{R}^M$
\begin{equation}
    X_{\theta} (s) = x_i \quad w.p. \quad p_i(s) := f_{\theta}^i(s)
\end{equation}

The DCOP update $Y_{\hat{\gamma}}X_\theta$ and this parametrization $X_\theta$ almost always would have disjoint supports. This could be tackled in practice projecting the sample DCOP update $\hat{Y}_{\hat{\gamma}}X_\theta$ onto the support of $X_\theta$ (i.e. given a sample transition $(s,a,s')$, we compute the DCOP update $\hat{Y}_{\hat{\gamma}}x_j =  \frac{\pi(a|s)}{\mu(a|s)}x_j + (1-\hat{\gamma})$ for each atom $x_j$, then distribute its probability $p_j(s)$ to the immediate neighbours of $\hat{Y}_{\hat{\gamma}}x_j$).

The corresponding pseudocode is detailed in Algorithm \ref{distCSalgorithm}, analogous to that of \cite{DRL}:

\setlength{\algomargin}{1.5em}
\SetAlCapHSkip{0em}
\begin{algorithm}[H]
\caption{Categorical CS Algorithm} \label{distCSalgorithm}
\DontPrintSemicolon
\SetKwComment{CustomComment}{\#}{}
\SetKwInput{KwIn}{\hspace{-1.5em} input}
\SetKwInput{KwOut}{\hspace{-1.5em} output}
\KwIn{A transition $s_{t-1},a_{t-1},s_t$}
    %$\frac{d_\pi}{d_\mu}(s_{t-1}) := \sum_i x_i p_i(s_{t-1})$\;
    $m_i = 0$, $i\in \{0,\dots, M-1\}$\;
    \For{$j\in \{0,\dots, M-1\}$}{
        \CustomComment{Compute the projection $\tilde{T}x_j$ onto the support $\{x_i\}$}
        $\tilde{T}x_j \leftarrow \left[ \hat{\gamma} \frac{\pi(a_{t-1}|s_{t-1})}{\mu(a_{t-1}|s_{t-1})} x_j + 1 - \hat{\gamma} \right]_{C_{min}}^{C_{max}}$\;
        $b_j \leftarrow (\tilde{T}x_j -C_{min}) / \Delta x \quad$ \CustomComment{note that $b_j \in [0, M-1]$}
        $l \leftarrow \lfloor b_j \rfloor$, $u \leftarrow \lceil b_j \rceil$\;
        \CustomComment{Distribute probability of $\tilde{T}x_j$}
        $m_l \leftarrow m_l + p_j(s_{t-1}) (u - b_j)$\;
        $m_u \leftarrow m_u + p_j(s_{t-1}) (b_j - l)$\;
    }
\KwOut{Cross-entropy loss $-\sum_i m_i \log(p_i(s_t))$}

\end{algorithm}

Observations:
\begin{itemize}
    \item $C_{min}, C_{max} \in \R$ are the predefined lower and upper value limits of the covariate shift ratio. In this case, $C_{min} \geq 0$.
    \item The support is evenly spaced in $[C_{min}, C_{max}]$; we consider the set of atoms $\{x_i = C_{min} + i \Delta x : 0 \leq i < M \}$, with $\Delta x = \frac{C_{max}-C_{min}}{M-1}$
    \item The \textit{behavioural policy} $\mu$ is simply the uniformly random policy, so 
    \begin{equation} \label{behaviourPolicy}
        \mu(a|s) = \frac{1}{|\mathcal{A}|} \quad \forall a \in \mathcal{A}
    \end{equation}
    for any state $s \in \mathcal{S}$.
    \item The \textit{target policy} $\pi$ is the $\epsilon$-greedy policy with respect to the estimated state-action values of the model. Hence,%expected value of the distribution output of the target network. Hence,
    \begin{equation} \label{targetPolicy}
        \pi_\theta (a|s) = 
        \left\{ 
            \begin{array}{lcc}
                (1-\epsilon)+\epsilon \frac{1}{|\mathcal{A}|} &   \text{if} & a = \text{arg max}_a Q_\theta(s,a) \\
                \epsilon \frac{1}{|\mathcal{A}|} &  \text{otherwise} &
            \end{array}
        \right. 
    \end{equation}
   for any state $s \in \mathcal{S}$.
   \item Considering equations \ref{behaviourPolicy} and \ref{targetPolicy}, we have that
   \begin{equation} \label{policyQuotient}
        \frac{\pi_\theta (a|s)}{\mu(a|s)} = 
        \left\{ 
            \begin{array}{lcc}
                |\mathcal{A}|(1-\epsilon)+\epsilon  &   \text{if} & a = \text{arg max}_a Q_\theta(s,a) \\
                \epsilon  &  \text{otherwise} &
            \end{array}
        \right. 
    \end{equation}
\end{itemize}


\subsection{Exponential Bins}

\subsection{Log-Categorical COP-TD}


\newpage 

\thispagestyle{plain}
\section{Evaluation of the proposal}



\newpage

\thispagestyle{plain}
\section{Conclusions}


\newpage

\appendix
\pagestyle{fancy}
\renewcommand{\sectionmark}[1]{\markboth{#1}{}}
\renewcommand{\subsectionmark}[1]{\markright{\thesubsection.\quad #1}}

\thispagestyle{plain}
\section*{Appendix}
\addcontentsline{toc}{section}{Appendix}
\sectionmark{Appendix}

\setcounter{section}{1}
\setcounter{theorem}{0}
\setcounter{definition}{0}
\setcounter{equation}{0}

\subsection{Mixture Distributions}

\begin{remark}
All random variables presented in this document are considered to be real-valued, i.e. their measurable space is $E= \R$. 
\end{remark}

A random variable $Y$ is a mixture distribution if it is derived from a collection of other random variables $\{X_i\}$, $i\in \{1, \dots, N\}$, (named mixture components) in such a way that the combination of these parent distributions is driven according to a certain distribution $A$ (called mixing distribution). $A$ encapsulates the mixture weights $\alpha_i \sim A$, $i\in \{1, \dots, N\}$, which represent the probabilities of each individual mixture component $X_i$. 

The mixture distribution $Y$ can be defined in terms of its density function $f_Y$, which is the resulting $\alpha$-convex combination of the mixture components' density functions:
\begin{equation} \label{mixtureDef}
    f_{Y} (x) = \sum_{i=1}^{N} \alpha _i f_{X_i}(x)
\end{equation}

Let's present some interesting properties of mixture distributions:

\begin{mixtureProperty} \label{mixtureProperty1}
	The expectation of the mixture distribution $Y$ is the convex combination of expectations of each mixture component: 
    \begin{equation} \label{mixtureProperty1eq}
    \begin{split}
        \E [Y] &= \int_{-\infty}^{\infty} x f_{Y}(x) dx = \int_{-\infty}^{\infty} x \sum_{i=1}^{N} \alpha_i f_{X_i}(x) dx \\
        &= \sum_{i=1}^{N} \alpha_i \int_{-\infty}^{\infty} x f_{X_i}(x) dx \\
        &= \sum_{i=1}^{N} \alpha_i \E [X_i]
    \end{split}
    \end{equation}
\end{mixtureProperty}


\begin{mixtureProperty} \label{mixtureProperty2}
	Let be $Z$ a mixture distribution with mixture components $\{g_i(X_i)\}$, $i\in \{1, \dots, N\}$ and mixing weights $\alpha_i \sim A$
    \begin{equation} \label{mixtureProperty2eq}
        \E [Z] = \sum_{i=1}^{N} \alpha_i \E [g_i(X_i)]
    \end{equation}
\end{mixtureProperty}


\begin{mixtureProperty} \label{mixtureProperty3}
	Let be $Z=g(Y)$, being $Y$ a mixture distribution with mixture components $\{X_i\}$, $i\in \{1, \dots, N\}$, and $g$ a monotonic, invertible and differentiable function. Then we have that $Z$ is a mixture distribution whose expectation is
    \begin{equation} \label{mixtureProperty3eq}
        \begin{split}
        \E [Z] &= \int_{-\infty}^{\infty} g(x) f_{Y}(x) dx \\
        &= \int_{-\infty}^{\infty} g(x) \left( \sum_{i=1}^{N} \alpha_i f_{X_i}(x)\right) dx = \sum_{i=1}^{N} \alpha_i \int_{-\infty}^{\infty} g(x) f_{X_i}(x) dx \\
        &= \sum_{i=1}^{N} \alpha_i \E [g(X_i)]
        \end{split}
    \end{equation}
    Note that in both the first and last steps the so-called Law of the Unconscious Statistician has been applied, which states that
    \begin{equation} \label{LOTUS}
    	\E [g(X)] = \int_{-\infty}^{\infty} y f_{g(X)}(y) dy  = \int_{-\infty}^{\infty} g(x) f_{X}(x) dx
    \end{equation}

\end{mixtureProperty}

We emphasize the relevance of \textbf{Property \ref{mixtureProperty3}}. In distributional TD learning, the distribution mixture plays the role of the expectation in expected TD. But while $\E[g(x)] \not= g(\E[x])$, we can interchange mixtures and functions. This allows us to circumvent Jensen's inequalities.

\begin{comment}
\subsubsection*{Sum of Distributions} 
	
The sum of two independent random variables $X_1$ and $X_2$, $Y=X_1 + X_2$, results in a random variable whose density function is the convolution of the density functions of each summand
\begin{equation}
	f_{Y}(y) = \int_{-\infty}^{\infty} f_{X_1}(y-x) f_{X_2}(x)dx = \left(f_{X_1}* f_{X_2} \right) (y)
\end{equation}

In the general case, considering a collection $\{X_i\}$, $i\in \{1, \dots, N\}$, of independent random variables, the density function of the random variable $Y = \sum_{i=1}^{N} X_i $
can be expressed as the convolution of all the individual density functions:
\begin{equation}
	f_{Y} (x) = \left(f_{X_1} * \cdots * f_{X_N} \right) (x)
\end{equation}

Convolutions are linear operators.
\end{comment}





\subsection{Metrics over Distributions} \label{metrics}


\subsubsection{Kullback-Leibler Divergence}

\begin{definition}
	Let be $\nu_1, \nu_2 \in \mathscr{P}(\R)$ two probability distributions. The Kullback-Leibler (KL) divergence of $\nu_1$ from $\nu_2$ is defined as
	\begin{equation*}
		D_{KL} (\nu_1,\nu_2) = \int_{-\infty}^{\infty} \nu_1(x) \log \frac{\nu_1(x)}{\nu_2(x)} dx
	\end{equation*}
\end{definition}

\subsubsection{Wasserstein}

\begin{definition}
	The Wasserstein distance $d_p$, for $p\in [1,\infty)$, between two distributions $\nu_1, \nu_2 \in \mathscr{P}(\R)$, with cumulative distribution functions $F_{\nu_1}, F_{\nu_2}$ respectively, can be defined by:
	\begin{equation*}
		d_p(\nu_1,\nu_2) = \left( \int_{\R} | F_{\nu_1}^{-1}(u) - F_{\nu_2}^{-1}(u)|^p du \right) ^{1/p}
	\end{equation*}
	Further, the supremum-Wasserstein metric $\bar{d}_p$ is defined between two value distribution functions $Z, Z' \in \Z$ by
	\begin{equation*}
		\bar{d}_p(Z,Z') = \sup_{(s,a) \in \Sspace \times \A} d_p(Z(s,a), Z'(s,a))
	\end{equation*}
\end{definition}


\subsubsection{Cramér}

\begin{definition}
	The Cramér distance $\ell_2$ between two distributions $\nu_1, \nu_2 \in \mathscr{P}(\R)$, with cumulative distribution functions $F_{\nu_1}, F_{\nu_2}$ respectively, is defined by:
	\begin{equation*}
		\ell_2(\nu_1,\nu_2) = \left( \int_{\R} (F_{\nu_1}(x) - F_{\nu_2}(x))^2dx \right) ^{1/2}
	\end{equation*}
	Further, the supremum-Cramér metric $\bar{\ell}_2$ is defined between two value distribution functions $Z, Z' \in \Z$ by
	\begin{equation*}
		\bar{\ell}_2(Z,Z') = \sup_{(s,a) \in \Sspace \times \A} \ell_2(Z(s,a), Z'(s,a))
	\end{equation*}
\end{definition}


%\subsubsection{Generalized Cramér}





\subsection{Contraction Mappings}



\subsection{Implementation Details}

TODO!!!

Our baseline is the C51 distributional reinforcement learning agent\cite{DRL} within Dopamine framework\cite{dopamine}. We use published hyperparameters unless otherwise noted. We augment the C51 network by adding an extra head, the distributional ratio model $X(s)$, to the final convolutional layer, whose role is to predict the distribution of the ratio $d_\pi / d_\mu$ . This model consists of a two-layer fully-connected network, with as many outputs as the number of atoms $M$ of the parametric model. A final softmax layer transforms the resulting logits into probabilities. 

\newpage

\thispagestyle{plain}
\addcontentsline{toc}{section}{Referencias}\sectionmark{Referencias}
\begin{thebibliography}{50}

	\bibitem{DRL} M. G. Bellamare, W. Dabney and R. Munos. A Distributional Perspective on Reinforcement Learning.
	\bibitem{COPTD} A. Hallak and S. Mannor. Consistent On-Line Off-Policy Evaluation.
	\bibitem{DCOPTD} C. Gelada and M. G. Bellamare. Off-Policy Deep Reinforcement Learning by Bootstrapping the Covariate Shift.
	\bibitem{unichain}  S. P. Meyn and R. L. Tweedie. Markov chains and stochastic stability. 2012.
	\bibitem{TD} Sutton, R. S., and Barto, A. G. 2018. Reinforcement learning: An introduction. MIT Press, 2nd edition.
	\bibitem{projectedBellman} Tsitsiklis, J. N., and Van Roy, B. 1997. An analysis of temporal-difference learning with function approximation. IEEE Transactions on Automatic Control 42(5):674–690.
	\bibitem{dopamine} Dopamine: {A} Research Framework for Deep Reinforcement Learning. \url{https://github.com/google/dopamine}
	\bibitem{rlformulation} Bertsekas, D. and Tsitsiklis, J. Neuro-Dynamic Programming. Athena Scientific, 1996.
	\bibitem{ETD} Sutton, R. S., Mahmood, A. R., and White, M. An emphatic approach to the problem of off-policy temporal- difference learning. arXiv:1503.04269, 2015.
	\bibitem{generalizedETD} Hallak, Assaf, Tamar, Aviv, Munos, Remi, and Mannor, Shie. Generalized emphatic temporal difference learning: Bias-variance analysis. arXiv preprint arXiv:1509.05172, 2015.
	\bibitem{DRLlinear} Marc G. Bellemare, Nicolas Le Roux, Pablo Samuel Castro and Subhodeep Moitra. Distributional reinforcement learning with linear function approximation
	\bibitem{analysisC51} Rowland, Mark, Bellemare, Marc G, Dabney, Will, Munos, Rémi, and Teh, Yee Whye. An analysis of categorical distributional reinforcement learning. In Proceedings of the International Conference on Artificial Intelligence and Statistics, 2018.
	\bibitem{firstDistributional} Jaquette, Stratton C. Markov decision processes with a new opti- mality criterion: Discrete time. The Annals of Statistics, 1(3): 496–505, 1973.
	\bibitem{ALE} Bellemare, Marc G, Naddaf, Yavar, Veness, Joel, and Bowling, Michael. The arcade learning environment: An evaluation plat- form for general agents. Journal of Artificial Intelligence Re- search, 47:253–279, 2013.





\end{thebibliography}


%\thispagestyle{plain}\nocite{*}\addcontentsline{toc}{section}{Referencias}\sectionmark{Referencias}\bibliography{Bibliografia} \bibliographystyle{unsrt}

\end{document}